\documentclass[11pt,reqno]{amsart}

% ==================================================
% 1. Packages
% ==================================================

\usepackage{amsmath,amssymb,amsthm,mathtools}
\usepackage{mathrsfs}
\usepackage{graphicx}
\usepackage{enumitem}
\usepackage{tikz}
\usepackage[hidelinks]{hyperref}
\usepackage[nameinlink,noabbrev]{cleveref}
% ==================================================
% 2. Mathematical operators
% ==================================================

\DeclareMathOperator{\Aut}{Aut}
\DeclareMathOperator{\Stab}{Stab}
\DeclareMathOperator{\Fix}{Fix}
\DeclareMathOperator{\rank}{rank}
\DeclareMathOperator{\im}{im}
\DeclareMathOperator{\Ker}{Ker}
\DeclareMathOperator{\PLC}{PLC}

\DeclareMathOperator{\Inv}{Inv}
% ==================================================
% 3. Mathematical symbols
% ==================================================

\newcommand{\R}{\mathbb{R}}
\newcommand{\C}{\mathbb{C}}
\newcommand{\Z}{\mathbb{Z}}
\newcommand{\N}{\mathbb{N}}

\newcommand{\norm}[1]{\left\lVert #1 \right\rVert}
\newcommand{\abs}[1]{\left\lvert #1 \right\rvert}
\newcommand{\pairing}[2]{\left\langle #1,#2 \right\rangle}

% Coxeter group notation (optional)

\newcommand{\simpleRoots}{\Pi}
\newcommand{\limitRoots}{\mathcal{E}}
\newcommand{\imaginaryCone}{\mathcal{I}}
\newcommand{\TitsCone}{\mathcal{T}}
\newcommand{\depth}{\operatorname{dp}}

% ==================================================
% 4. Theorem environments
% ==================================================

\theoremstyle{plain}

\newtheorem{theorem}{Theorem}[section]
\newtheorem{lemma}[theorem]{Lemma}
\newtheorem{proposition}[theorem]{Proposition}
\newtheorem{corollary}[theorem]{Corollary}

\theoremstyle{definition}

\newtheorem{definition}[theorem]{Definition}
\newtheorem{example}[theorem]{Example}
\newtheorem{notation}[theorem]{Notation}

\theoremstyle{remark}

\newtheorem{remark}[theorem]{Remark}

% ==================================================
% 5. Equation numbering
% ==================================================

\numberwithin{equation}{section}

% ==================================================
% 6. Hyper-reference names
% ==================================================

\crefname{theorem}{Theorem}{Theorems}
\crefname{lemma}{Lemma}{Lemmas}
\crefname{proposition}{Proposition}{Propositions}
\crefname{corollary}{Corollary}{Corollaries}
\crefname{definition}{Definition}{Definitions}
\crefname{equation}{Equation}{Equations}
\crefname{section}{Section}{Sections}
\DeclareMathOperator{\tr}{tr}
\DeclareMathOperator{\supp}{supp}
\DeclareMathOperator{\Pos}{Pos}
\DeclareMathOperator{\rad}{rad}
\DeclareMathOperator{\dom}{dom}

% ==================================================
% 7. Title and author
% ==================================================

\title[Some open problems of Dyer]
{Some open problems of Dyer}

\author{Haoquan Wang}

\address{
School of Mathematical Science,
Peking University,
Peking, China
}

\email{2300010755@stu.pku.edu.cn}

% Optional: AMS classification and keywords


\keywords{
Coxeter groups, Fuchsian groups, limit roots,
dominance order
}

% ==================================================
% 8. Document
% ==================================================

\setlength{\emergencystretch}{2em}
\begin{document}

% Abstract

\begin{abstract}

\end{abstract}

\maketitle
The following questions are organized by the eight themes in
\cite[Section~7]{DHR}. A check mark indicates a question resolved
in this manuscript, including a negative answer by counterexample.
Partial results are identified separately.

\par\medskip
\noindent\textbf{Literature scope (16 November 2026).}
The notes below distinguish published results, preprints, elementary
consequences, and the results claimed in this manuscript. ``No general
resolution was verified'' reports the scope of this search, not a proof
that a problem is still open. Numbered results refer to the versions
identified in the bibliography. Existing manuscript proofs are retained;
the literature annotations do not constitute an independent validation
of every new theorem here.
\tableofcontents
\section{Geometric Characteristics of \texorpdfstring{$E(\Phi)$}{E(Phi)} and \texorpdfstring{$\mathcal{Z}(\Phi)$}{Z(Phi)}}

\subsection{Intersection of the isotropic cone with the convex hull}

An important open question is whether the equality
\[
E(\Phi)=\widehat{\mathcal{Q}}\cap\operatorname{conv}(E(\Phi))
\]
holds for general irreducible Coxeter systems.
This equality has already been established in the weakly
hyperbolic case by \cite[Theorem~4.10]{DHR}.
An affirmative answer would provide a natural fractal
description of the limit-root set $E(\Phi)$.

The following counterexample shows that this equality fails in general.


\par\medskip
\noindent\textbf{Status.} \(\checkmark\) Answered negatively in this manuscript. The rank-five counterexample below disproves the equality in general; the Lorentzian case is known. This status depends on verification of the manuscript's counterexample.

\par\medskip
\noindent\textbf{Previous work and verified progress.}
Dyer--Hohlweg--Ripoll prove the weakly hyperbolic case in
\cite[Theorem~4.10]{DHR}. Higashitani--Mineyama--Nakashima independently
obtain the Lorentzian fractal description \cite[Theorem~1.2]{HMN2019}.
Neither statement covers arbitrary signature. The negative answer below
is a result claimed and proved in this manuscript; it is not being
attributed to those references.

\begin{proposition}\label{prop:counterexample-convE-Q}
There exists an irreducible based root system of rank five and
signature $(3,2)$ for which
\[
E \subsetneq \operatorname{conv}(E)\cap Q.
\]
More precisely, let
\[
V=\bigoplus_{i=1}^{5}\mathbb{R}\alpha_i
\]
and let $B$ be the symmetric bilinear form whose Gram matrix,
with respect to the basis $(\alpha_1,\ldots,\alpha_5)$, is
\[
G=
\begin{pmatrix}
1&-2&0&0&-2\\
-2&1&0&0&-2\\
0&0&1&-2&-2\\
0&0&-2&1&-2\\
-2&-2&-2&-2&1
\end{pmatrix}.
\]
Let $W=\langle s_1,\ldots,s_5\rangle$, where
\[
s_i(v)=v-2B(v,\alpha_i)\alpha_i,
\]
and let
\[
\Phi=W\{\alpha_1,\ldots,\alpha_5\},
\qquad
P=\sum_{i=1}^{5}\mathbb{R}_{\geq 0}\alpha_i,
\qquad
\Phi^+=\Phi\cap P.
\]
Define
\[
h\left(\sum_{i=1}^{5}c_i\alpha_i\right)=\sum_{i=1}^{5}c_i,
\qquad
\widehat{v}=\frac{v}{h(v)},
\qquad
E=\operatorname{Acc}\{\widehat{\rho}:\rho\in\Phi^+\},
\]
and put
\[
Q=\{v\in V:B(v,v)=0\}.
\]
Set
\[
\tau=2-\sqrt{3},
\qquad
x=\frac{\alpha_1+\tau\alpha_2}{1+\tau},
\qquad
y=\frac{\alpha_3+\tau\alpha_4}{1+\tau}.
\]
Then
\[
x,y\in E,
\qquad
B(x,y)=0,
\qquad
\frac{x+2y}{3}\notin E.
\]
\end{proposition}

\begin{proof}
The off-diagonal entries of $G$ belong to $\{0,-2\}$, so $G$
defines a based root system: an entry $0$ corresponds to
$m_{ij}=2$, whereas an entry $-2$ corresponds to $m_{ij}=\infty$.
Its Coxeter graph is connected, since the fifth vertex is
joined to each of the other four vertices. Hence the root
system is irreducible.

Write
\[
x_+=x,
\qquad
x_-=\frac{\tau\alpha_1+\alpha_2}{1+\tau},
\]
and
\[
y_+=y,
\qquad
y_-=\frac{\tau\alpha_3+\alpha_4}{1+\tau}.
\]
These are the limit roots of the two non-affine dihedral
root subsystems supported on
$\{\alpha_1,\alpha_2\}$ and $\{\alpha_3,\alpha_4\}$.

For completeness, the matrix of $s_1s_2$ on
$\operatorname{span}\{\alpha_1,\alpha_2\}$ is
\[
\begin{pmatrix}
15&-4\\
4&-1
\end{pmatrix}.
\]
Its eigenvectors $x_+$ and $x_-$ have respective eigenvalues
\[
\tau^{-2}=7+4\sqrt{3},
\qquad
\tau^2=7-4\sqrt{3}.
\]
Consequently,
\[
\widehat{(s_1s_2)^n\alpha_1}\longrightarrow x_+.
\]
The corresponding argument for $s_3s_4$ gives
\[
\widehat{(s_3s_4)^n\alpha_3}\longrightarrow y_+.
\]
Thus $x_+,y_+\in E$.

Direct calculation gives
\[
B(x_+,x_+)=B(x_-,x_-)
=B(y_+,y_+)=B(y_-,y_-)=0,
\]
\[
B(x_+,x_-)=B(y_+,y_-)=-1,
\]
and
\[
B(x_\varepsilon,y_\delta)=0
\qquad
(\varepsilon,\delta\in\{+,-\}).
\]
Let
\[
\eta=\alpha_5-2\alpha_1-2\alpha_2-2\alpha_3-2\alpha_4.
\]
Then
\[
\eta\perp_B
\operatorname{span}\{x_+,x_-,y_+,y_-\},
\qquad
B(\eta,\eta)=17.
\]
The vectors $(x_+,x_-,y_+,y_-,\eta)$ form a basis of $V$,
and the Gram matrix in this basis is
\[
\begin{pmatrix}
0&-1&0&0&0\\
-1&0&0&0&0\\
0&0&0&-1&0\\
0&0&-1&0&0\\
0&0&0&0&17
\end{pmatrix}.
\]
Therefore $B$ has signature $(3,2)$.

For $v=\sum_{i=1}^{5}c_i\alpha_i$, define linear forms
\[
\begin{aligned}
A_+(v)
&=\frac{1+\sqrt{3}}{2}c_1
  -\frac{\sqrt{3}-1}{2}c_2+2c_5,\\
A_-(v)
&=-\frac{\sqrt{3}-1}{2}c_1
  +\frac{1+\sqrt{3}}{2}c_2+2c_5,\\
C_+(v)
&=\frac{1+\sqrt{3}}{2}c_3
  -\frac{\sqrt{3}-1}{2}c_4+2c_5,\\
C_-(v)
&=-\frac{\sqrt{3}-1}{2}c_3
  +\frac{1+\sqrt{3}}{2}c_4+2c_5.
\end{aligned}
\]
These are the first four coordinates in the preceding basis:
\[
v=
A_+(v)x_++A_-(v)x_-
+C_+(v)y_++C_-(v)y_-+c_5\eta.
\]

We claim that every positive root $\rho$ satisfies
\begin{equation}\label{eq:root-gap-inequality}
\bigl(20A_+(\rho)-9C_+(\rho)\bigr)
\bigl(3A_+(\rho)-2C_+(\rho)\bigr)\geq 0.
\end{equation}
We first establish a useful description of positive roots.

Since the simple reflections have integer matrices in the
simple-root basis, every root has integer coordinates.
Let $\rho=\sum_i c_i\alpha_i$ be a positive root that is
not simple. Since
\[
1=B(\rho,\rho)=\sum_i c_i B(\rho,\alpha_i),
\]
there exists $i$ with $B(\rho,\alpha_i)>0$.
A simple reflection preserves all positive roots other than
its corresponding simple root. Hence
\[
s_i\rho\in\Phi^+,
\qquad
h(s_i\rho)=h(\rho)-2B(\rho,\alpha_i)<h(\rho).
\]
The height is a positive integer, so repeating this procedure
eventually reaches a simple root.

Reversing the resulting sequence shows that every positive
root can be obtained from a simple root by a sequence of
simple reflections in which each step strictly increases
one simple-root coordinate and leaves the others unchanged.

Suppose first that the fifth coordinate of $\rho$ is positive.
Choose such a coordinate-increasing sequence.
If the sequence contains $s_5$, let $v$ be the root immediately
after its last occurrence.
If it contains no $s_5$, then it must start at $\alpha_5$;
in this case put $v=\alpha_5$.
In either case,
\[
\rho=uv
\qquad\text{for some}\qquad
u\in W_{\{1,2,3,4\}}.
\]

Write
\[
v=\sum_{i=1}^{4}c_i\alpha_i+d\alpha_5,
\qquad
S=c_1+c_2+c_3+c_4.
\]
We have
\begin{equation}\label{eq:seed-coordinate-bound}
d\geq 2S.
\end{equation}
Indeed, if $v=\alpha_5$, this is immediate.
Otherwise, let $d_0$ be the fifth coordinate immediately
before the last occurrence of $s_5$.
Since that step increases the fifth coordinate,
\[
d_0-2S
=
B\left(\sum_{i=1}^{4}c_i\alpha_i+d_0\alpha_5,\alpha_5\right)
<0.
\]
The fifth coordinate after reflection is
\[
d=4S-d_0>2S,
\]
which proves \eqref{eq:seed-coordinate-bound}.

Since $c_i\geq 0$ and $S\leq d/2$, the definitions of the
four linear forms imply
\[
\frac{9-\sqrt{3}}{4}\,d
\leq
A_\varepsilon(v),C_\varepsilon(v)
\leq
\frac{9+\sqrt{3}}{4}\,d
\qquad
(\varepsilon\in\{+,-\}).
\]
In particular, these four values are strictly positive.
Put
\[
R=\frac{9+\sqrt{3}}{9-\sqrt{3}}.
\]
Then
\begin{equation}\label{eq:seed-ratio-bound}
R^{-1}
\leq
\frac{A_\varepsilon(v)}{C_\delta(v)}
\leq R
\qquad
(\varepsilon,\delta\in\{+,-\}).
\end{equation}
Moreover,
\begin{equation}\label{eq:elementary-gap-bounds}
R<\frac32,
\qquad
0<\tau<\frac3{10}.
\end{equation}

The action of the first two simple reflections on the
coordinates $(A_+,A_-)$ is
\[
\begin{pmatrix}
A_+(s_1v)\\
A_-(s_1v)
\end{pmatrix}
=
\begin{pmatrix}
0&\tau^{-1}\\
\tau&0
\end{pmatrix}
\begin{pmatrix}
A_+(v)\\
A_-(v)
\end{pmatrix},
\]
and
\[
\begin{pmatrix}
A_+(s_2v)\\
A_-(s_2v)
\end{pmatrix}
=
\begin{pmatrix}
0&\tau\\
\tau^{-1}&0
\end{pmatrix}
\begin{pmatrix}
A_+(v)\\
A_-(v)
\end{pmatrix}.
\]
They fix the coordinates $(C_+,C_-)$ and the vector $\eta$.
The reflections $s_3,s_4$ act by the same two matrices on
$(C_+,C_-)$ and fix $(A_+,A_-)$ and $\eta$.

It follows that there exist integers $k,\ell$ and signs
$\varepsilon,\delta\in\{+,-\}$ such that
\[
A_+(\rho)=\tau^k A_\varepsilon(v),
\qquad
C_+(\rho)=\tau^\ell C_\delta(v).
\]
In particular, $A_+(\rho),C_+(\rho)>0$, and
\[
\frac{A_+(\rho)}{C_+(\rho)}
\in
\tau^{k-\ell}[R^{-1},R].
\]
If $k-\ell\leq 0$, then
\[
\frac{A_+(\rho)}{C_+(\rho)}
\geq R^{-1}>\frac23.
\]
If $k-\ell\geq 1$, then
\[
\frac{A_+(\rho)}{C_+(\rho)}
\leq \tau R<\frac9{20}.
\]
Thus the ratio cannot lie between $9/20$ and $2/3$.
This proves \eqref{eq:root-gap-inequality} whenever the
fifth coordinate of $\rho$ is positive.

Suppose now that the fifth coordinate of $\rho$ is zero.
Along a coordinate-increasing sequence producing $\rho$,
the fifth coordinate must remain zero. Hence the sequence
does not use $s_5$ and starts at one of
$\alpha_1,\ldots,\alpha_4$.
Because the two blocks $\{1,2\}$ and $\{3,4\}$ are orthogonal,
the resulting root is supported entirely on one of them.

If $\rho$ is supported on $\{\alpha_1,\alpha_2\}$, then
$C_+(\rho)=0$, and the left-hand side of
\eqref{eq:root-gap-inequality} is
\[
60A_+(\rho)^2\geq 0.
\]
If $\rho$ is supported on $\{\alpha_3,\alpha_4\}$, then
$A_+(\rho)=0$, and the same expression equals
\[
18C_+(\rho)^2\geq 0.
\]
This completes the proof of
\eqref{eq:root-gap-inequality} for every positive root.

Define the continuous homogeneous quadratic function
\[
F(v)=
\bigl(20A_+(v)-9C_+(v)\bigr)
\bigl(3A_+(v)-2C_+(v)\bigr).
\]
For every positive root $\rho$,
\[
F(\widehat{\rho})=\frac{F(\rho)}{h(\rho)^2}\geq 0.
\]
Passing to accumulation points gives
\[
F(e)\geq 0
\qquad\text{for every }e\in E.
\]

Finally, put
\[
z=\frac{x_++2y_+}{3}.
\]
Since $x_+,y_+\in E$ are isotropic and orthogonal,
\[
z\in\operatorname{conv}(E)\cap Q.
\]
On the other hand, its coordinates in the chosen basis give
\[
A_+(z)=\frac13,
\qquad
C_+(z)=\frac23.
\]
Therefore
\[
F(z)
=
\left(\frac{20}{3}-6\right)
\left(1-\frac43\right)
=
-\frac29<0.
\]
Consequently $z\notin E$, proving
\[
E\subsetneq\operatorname{conv}(E)\cap Q.
\]
\end{proof}

\subsection{Homeomorphism types of limit-root sets}

Another natural problem is to characterize pairs of root
systems $(\Phi,\Phi')$ for which
\[
E(\Phi)\cong E(\Phi')
\]
as topological spaces.



\par\medskip
\noindent\textbf{Status.} Partially solved. Boundary models are known for several classes; no classification in the full stated generality was verified.

\par\medskip
\noindent\textbf{Previous work and verified progress.}
In the Lorentzian setting, the identification with a hyperbolic limit
set \cite{HPR2020} makes boundary methods available. Mineyama constructs
an equivariant continuous surjection onto the limit set
\cite[Theorem~1.1]{Mineyama2014}; his boundary for nonhyperbolic groups
uses an extended definition, and the map is not asserted to be a
homeomorphism. A general classification of pairs of realizations with
homeomorphic limit-root sets was not located in the sources checked.

\subsection{Recovering the root system from its limit roots}

More generally, one may ask to what extent the set of
limit roots $E(\Phi)$ determines the root system $\Phi$.


\par\medskip
\noindent\textbf{Status.} Partially solved. Topology alone does not recover the realization, as non-affine dihedral examples show. Reconstruction from embedded or marked data is a separate question.

\par\medskip
\noindent\textbf{Previous work and verified progress.}
The rank-two computations in \cite[Section~2.3]{HLR} already show
that topology alone cannot recover a realization: all non-affine dihedral
realizations have a two-point limit set, while their Gram parameter can
vary. Conversely, changing an infinite-edge parameter from $-1$ to a value
less than $-1$ changes a singleton limit set into a two-point set.
Thus neither the abstract Coxeter presentation nor the unmarked topology
of $E$ determines all the other data. Recovery from an embedded set or
from a marked $W$-space is a different question.

\subsection{Connectedness}
Which root systems have a connected limit-root set $E(\Phi)$?


\par\medskip
\noindent\textbf{Status.} Partially solved. Particular classes are understood; no general connectedness criterion was verified.

\par\medskip
\noindent\textbf{Previous work and verified progress.}
For a Lorentzian realization with the entire normalized isotropic
sphere contained in the simple-root simplex, \cite[Theorem~1.2(a)]{HMN2019}
gives $E=\widehat Q$; this is connected in rank at least three.
In settings where $E$ is identified homeomorphically with a group boundary,
boundary connectedness criteria also apply. Such an identification must
be established for the particular representation; it does not follow
from the signature alone. No general root-system classification was
verified in this search.

\subsection{Local connectedness}
Which root systems have a locally connected limit-root set $E(\Phi)$?


\par\medskip
\noindent\textbf{Status.} Partially solved. Boundary local-connectedness results apply when the required identification with the limit-root set is established.

\par\medskip
\noindent\textbf{Previous work and verified progress.}
Dasgupta--Hruska prove local connectedness of the Bowditch boundary
for relatively one-ended relatively hyperbolic pairs
\cite{DH2024}, removing earlier restrictions on peripheral groups.
This settles the corresponding case for $E$ whenever an equivariant
homeomorphism from that boundary to $E$ is available. It is a conditional
application, not a theorem that every limit-root set is locally connected;
a continuous boundary surjection alone should not be described as a
homeomorphism.

\subsection{Total disconnectedness}
Which root systems have a totally disconnected limit-root set $E(\Phi)$?


\par\medskip
\noindent\textbf{Status.} Partially solved. Finite and Cantor-type examples are known; no general characterization was verified.

\par\medskip
\noindent\textbf{Previous work and verified progress.}
The dihedral examples above provide finite totally disconnected
sets. In a convex-cocompact Lorentzian realization of a non-elementary
virtually free group, the hyperbolic limit set, hence $E$ by
\cite{HPR2020}, is a Cantor set. This is a consequence of the boundary
model, not a classification in arbitrary signature. In particular,
abstract virtual freeness does not suffice without controlling the
realization and possible parabolic elements.

\subsection{Exposed points and extreme points}

Let
\[
C=\operatorname{conv}(E(\Phi)),
\]
and denote its sets of exposed points and extreme points by
$E_{\mathrm{exp}}(\Phi)$ and $E_{\mathrm{ext}}(\Phi)$,
respectively.
Is the equality
\[
E_{\mathrm{exp}}(\Phi)=E_{\mathrm{ext}}(\Phi)
\]
always valid?

\par\medskip
\noindent\textbf{Status.}
\(\checkmark\) \textit{Answered negatively for general based root systems.}
The rank-eight construction below has signature $(5,3)$ and
provides two extreme points of $\operatorname{conv}(E)$
which are not exposed.
The construction uses the Gram parameter $-2$ on infinite edges.
It does not settle the question for the standard Tits realization,
where the parameter on every infinite edge is fixed at $-1$.

\par\medskip
\noindent\textbf{Previous work and verified progress.}
Dyer--Hohlweg--Ripoll discuss extreme and exposed limit roots in
\cite[Section~3.3]{DHR} and establish the density statement
\[
\overline{E_{\mathrm{ext}}(\Phi)}=E(\Phi)
\]
under the hypotheses of \cite[Theorem~3.4]{DHR}.
This does not imply that every extreme point is exposed.
The following construction gives a counterexample to that
stronger assertion.

\begin{proposition}\label{prop:nonexposed-extreme-limit-roots}
There exists an irreducible based root system $\Phi$ of rank eight
and signature $(5,3)$ such that
\[
E_{\mathrm{exp}}(\Phi)
\subsetneq E_{\mathrm{ext}}(\Phi).
\]
More precisely, its limit-root convex hull has two distinct
extreme points $x,y$ such that every supporting hyperplane at
either point contains both points.
\end{proposition}

\begin{proof}
Let $V=\bigoplus_{i=1}^{8}\mathbb R\alpha_i$,
and define $B$ by the Gram matrix
\[
G=
\begin{pmatrix}
1&-2&0&0&0&0&0&0\\
-2&1&-2&0&0&0&0&0\\
0&-2&1&-2&0&0&0&0\\
0&0&-2&1&-2&0&0&0\\
0&0&0&-2&1&-2&0&0\\
0&0&0&0&-2&1&-2&0\\
0&0&0&0&0&-2&1&-2\\
0&0&0&0&0&0&-2&1
\end{pmatrix}.
\]
Let
\[
s_i(v)=v-2B(v,\alpha_i)\alpha_i,
\qquad
W=\langle s_1,\ldots,s_8\rangle,
\qquad
\Phi=W\{\alpha_1,\ldots,\alpha_8\}.
\]
These data define a based root system, and its Coxeter graph
is connected.
The eigenvalues of $G$ are
\[
1-4\cos\frac{k\pi}{9},
\qquad 1\le k\le8.
\]
Exactly three are negative, so the signature is $(5,3)$.

Use the normalization
\[
h\left(\sum_i c_i\alpha_i\right)=\sum_i c_i,
\qquad
\widehat v=\frac{v}{h(v)}.
\]
Write
\[
E=E(\Phi),
\qquad
C=\operatorname{conv}(E),
\qquad
K=\mathbb R_{\ge0}C.
\]
The cone $K$ is closed and $W$-invariant, is contained in the
nonnegative simple-root cone, and satisfies
\[
B(v,w)\le0
\qquad(v,w\in K).
\]
These are standard properties of the closed imaginary cone;
see \cite[Section~2]{DHR}.

\medskip
\noindent\textit{The three dihedral blocks.}
Put
\[
(a_1,b_1)=(1,2),
\qquad
(a_2,b_2)=(4,5),
\qquad
(a_3,b_3)=(7,8),
\]
and set
\[
\tau=2-\sqrt3,
\qquad
\lambda=\tau^{-2}=7+4\sqrt3.
\]
For $1\le i\le3$, define
\[
\xi_i=
\frac{\alpha_{a_i}+\tau\alpha_{b_i}}{1+\tau},
\qquad
\eta_i=
\frac{\tau\alpha_{a_i}+\alpha_{b_i}}{1+\tau}.
\]
These are the two limit roots of the corresponding non-affine
dihedral subsystem. In particular,
\[
\xi_i,\eta_i\in E,
\qquad
B(\xi_i,\eta_i)=-1.
\]
The three two-dimensional blocks are pairwise orthogonal.

We will use that $\xi_1$ and $\xi_3$ generate extreme rays of $K$.
Indeed, for $i\in\{1,3\}$, the functional
\[
v\longmapsto
B(\xi_i,v)-\sum_{j\notin\{a_i,b_i\}}c_j(v)
\]
is nonpositive on $C$.
Equality forces $v$ to lie in the corresponding dihedral block
and to be orthogonal to $\xi_i$.
Since that block is nondegenerate of signature $(1,1)$,
its intersection with $\xi_i^\perp$ is $\mathbb R\xi_i$.
The normalization $h(v)=1$ then gives $v=\xi_i$.
Thus $\xi_1$ and $\xi_3$ are exposed points of $C$.

Define
\[
a=s_1s_2,
\qquad
b=s_4s_5,
\qquad
c=s_7s_8,
\qquad
T=s_3s_6.
\]
The elements $a,b,c$ commute. On their respective blocks,
they have eigenvalues $\lambda,\lambda^{-1}$ with
eigenvectors $\xi_i,\eta_i$.

\medskip
\noindent\textit{A family approaching a rank-two projection.}
Define
\[
P_i^+(v)=-B(v,\eta_i)\xi_i,
\qquad
P_i^-(v)=-B(v,\xi_i)\eta_i,
\]
and
\[
P_U=I-\sum_{i=1}^{3}(P_i^++P_i^-).
\]
Here $P_U$ is the projection onto the orthogonal complement of
the sum of the three dihedral blocks.

Put $\gamma=7/3$, and define
\[
\begin{aligned}
A_0&=\gamma^{-1}(P_1^++P_3^+)T,\\
A_1&=\gamma^{-1}P_2^+T,\\
A_2&=\gamma^{-1}P_UT,\\
A_3&=\gamma^{-1}P_2^-T,\\
A_4&=\gamma^{-1}(P_1^-+P_3^-)T.
\end{aligned}
\]
For every integer $d\ge1$, set
\[
g(d)=a^{2d}b^dc^{2d}T,
\qquad
F(d)=\frac{g(d)}{\gamma\lambda^{2d}}.
\]
The spectral decomposition of the three commuting blocks gives
the exact identity
\begin{equation}\label{eq:nonexposed-factor-expansion}
F(d)=
A_0+\lambda^{-d}A_1+\lambda^{-2d}A_2
+\lambda^{-3d}A_3+\lambda^{-4d}A_4.
\end{equation}

Direct calculation gives
\begin{equation}\label{eq:nonexposed-projection}
A_0^2=A_0,
\qquad
\operatorname{im}A_0=\operatorname{span}\{\xi_1,\xi_3\},
\qquad
A_0\xi_1=\xi_1,
\qquad
A_0\xi_3=\xi_3.
\end{equation}
Moreover, with
\[
u=\frac{8+4\sqrt3}{7},
\qquad
v=\frac{8-4\sqrt3}{7},
\]
we have $u,v>0$ and
\begin{equation}\label{eq:nonexposed-transitions}
A_0\xi_2=u\xi_1+v\xi_3,
\qquad
A_1\xi_1=v\xi_2,
\qquad
A_1\xi_3=u\xi_2.
\end{equation}
For example, the matrix of $(P_1^++P_2^++P_3^+)T$
on the ordered vectors $(\xi_1,\xi_2,\xi_3)$ is
\[
\begin{pmatrix}
\frac73&\frac{8+4\sqrt3}{3}&0\\
\frac{8-4\sqrt3}{3}&\frac{11}{3}&\frac{8+4\sqrt3}{3}\\
0&\frac{8-4\sqrt3}{3}&\frac73
\end{pmatrix},
\]
from which these identities follow.

\medskip
\noindent\textit{Convergence of the products.}
Set
\[
d_n=n+10,
\qquad
\varepsilon_n=\lambda^{-d_n},
\qquad
F_n=F(d_n),
\qquad
L_n=F_1\cdots F_n.
\]
Each $L_n$ is a positive scalar multiple of a group element,
and hence
\[
L_nK=K,
\qquad
L_n^{-1}K=K.
\]

For completeness, we verify convergence in operator norm.
Equip $V$ with the norm
\[
\|z\|_*=
\|A_0z\|_\infty+\|(I-A_0)z\|_\infty,
\]
where $\|\cdot\|_\infty$ is the maximum norm in the simple-root
coordinates. Since $A_0$ is a projection,
\[
\|A_0\|_*\le1,
\qquad
\|I-A_0\|_*\le1.
\]
Write
\[
F_n=A_0+R_n,
\qquad
\delta_n=\|R_n\|_*.
\]

The explicit matrices above give the elementary bounds
\[
\|A_0\|_\infty<6,
\qquad
\|A_j\|_\infty<8
\quad(1\le j\le4).
\]
Consequently,
\[
\|z\|_\infty\le\|z\|_*\le13\|z\|_\infty,
\]
and \eqref{eq:nonexposed-factor-expansion} implies, with a
deliberately coarse bound,
\[
\delta_n\le5000\varepsilon_n.
\]
In particular,
\begin{equation}\label{eq:nonexposed-small-perturbation}
S:=\sum_{n\ge1}\delta_n
\le
\frac{5000\lambda^{-11}}{1-\lambda^{-1}}
<10^{-8}.
\end{equation}

We have the uniform bound
\[
\|L_n\|_*
\le
\prod_{j=1}^{n}(1+\delta_j)
\le e^S.
\]
Also,
\[
L_n(I-A_0)=L_{n-1}R_n(I-A_0),
\]
so
\[
\|L_n(I-A_0)\|_*\le e^S\delta_n.
\]
It follows that
\[
\begin{aligned}
L_{n+1}-L_n
&=L_nR_{n+1}-L_n(I-A_0),\\
\|L_{n+1}-L_n\|_*
&\le e^S(\delta_{n+1}+\delta_n).
\end{aligned}
\]
Thus $(L_n)$ converges to an operator $L$.
The preceding estimate also yields
\[
L(I-A_0)=0,
\qquad\text{hence}\qquad
L=LA_0.
\]

Comparing the products $F_1\cdots F_n$ and $A_0^n=A_0$
by telescoping gives
\[
\|L-A_0\|_*\le e^S-1<1.
\]
Therefore $L$ is injective on $\operatorname{im}A_0$.
Together with $L=LA_0$, this proves
\[
\operatorname{rank}L=2.
\]

Applying the same argument to each tail gives operators
\[
Q_n=\lim_{m\to\infty}F_{n+1}\cdots F_m
\]
such that
\begin{equation}\label{eq:nonexposed-tail}
L=L_nQ_n,
\qquad
Q_n\longrightarrow A_0.
\end{equation}

\medskip
\noindent\textit{The two limiting extreme rays.}
Put
\[
X=L\xi_1,
\qquad
Y=L\xi_3,
\qquad
x=\widehat X,
\qquad
y=\widehat Y.
\]
The vectors $X,Y$ belong to $K$, are nonzero, and are linearly
independent. In particular, $h(X),h(Y)>0$.
Since normalized group actions preserve $E$,
\[
\widehat{L_n\xi_1},\widehat{L_n\xi_3}\in E.
\]
Passing to the limit gives $x,y\in E$, with $x\ne y$.

We show that $X$ generates an extreme ray of $K$.
Suppose
\[
X=A+B,
\qquad A,B\in K.
\]
For every $n$, put
\[
A_n=L_n^{-1}A,
\qquad
B_n=L_n^{-1}B.
\]
These vectors belong to $K$, and
\[
A_n+B_n=L_n^{-1}X=Q_n\xi_1\longrightarrow\xi_1.
\]
Because $K$ is contained in the nonnegative coordinate cone,
both sequences $(A_n)$ and $(B_n)$ are bounded.
After passing to a subsequence, write
\[
A_n\longrightarrow A_\infty,
\qquad
B_n\longrightarrow B_\infty.
\]
Then
\[
A_\infty,B_\infty\in K,
\qquad
A_\infty+B_\infty=\xi_1.
\]
Since $\xi_1$ generates an extreme ray,
$A_\infty=t\xi_1$ for some $t\ge0$.
Therefore
\[
A=L_nA_n\longrightarrow L(t\xi_1)=tX.
\]
Thus every decomposition of $X$ in $K$ is along its ray.
The same argument applies to $Y$, and hence
\begin{equation}\label{eq:nonexposed-extremality}
x,y\in E_{\mathrm{ext}}(\Phi).
\end{equation}

\medskip
\noindent\textit{The supporting hyperplanes cannot separate the endpoints.}
Set
\[
X_n=L_n\xi_1,
\qquad
Y_n=L_n\xi_3,
\qquad
Z=L\xi_2.
\]
Since $L=LA_0$, equation
\eqref{eq:nonexposed-transitions} gives
\begin{equation}\label{eq:nonexposed-mixed-direction}
Z=uX+vY.
\end{equation}

Using \eqref{eq:nonexposed-factor-expansion} and
\eqref{eq:nonexposed-transitions}, together with the uniform
boundedness of $L_n$, we obtain
\[
X_{n+1}-X_n
=
v\varepsilon_{n+1}L_n\xi_2
+O(\varepsilon_{n+1}^2).
\]
Put
\[
t_n=\sum_{k>n}\varepsilon_k.
\]
Since $L_n\xi_2\to Z$ and
$\sum_{k>n}\varepsilon_k^2=o(t_n)$, summing the preceding
identity gives
\[
X-X_n=vt_nZ+o(t_n).
\]
By \eqref{eq:nonexposed-mixed-direction},
\[
X_n=(1-uvt_n)X-v^2t_nY+o(t_n).
\]
After normalization, this becomes
\begin{equation}\label{eq:nonexposed-tangent-x}
\widehat{X_n}
=
x+
v^2\frac{h(Y)}{h(X)}t_n(x-y)
+o(t_n).
\end{equation}

Let $\ell$ be any linear functional attaining its maximum
over $C$ at $x$.
As $\widehat{X_n}\in E\subseteq C$,
\[
\ell(\widehat{X_n})-\ell(x)\le0.
\]
Divide \eqref{eq:nonexposed-tangent-x}, after applying $\ell$,
by $t_n$ and let $n\to\infty$.
Because $v^2h(Y)/h(X)>0$, this gives
\[
\ell(x)\le\ell(y).
\]
The maximality at $x$ gives the reverse inequality, so
\[
\ell(x)=\ell(y).
\]
Since $x\ne y$, the point $x$ is not exposed.

Likewise,
\[
Y_{n+1}-Y_n
=
u\varepsilon_{n+1}L_n\xi_2
+O(\varepsilon_{n+1}^2),
\]
and therefore
\[
\widehat{Y_n}
=
y+
u^2\frac{h(X)}{h(Y)}t_n(y-x)
+o(t_n).
\]
The same argument shows that every linear functional attaining
its maximum at $y$ also attains that maximum at $x$.
Thus $y$ is not exposed either.

Combining this with \eqref{eq:nonexposed-extremality} yields
\[
x,y\in
E_{\mathrm{ext}}(\Phi)\setminus E_{\mathrm{exp}}(\Phi),
\]
as required.
\end{proof}

\begin{remark}
The counterexample points are specified by the group elements
\[
g_n=
(s_1s_2)^{2(n+10)}
(s_4s_5)^{n+10}
(s_7s_8)^{2(n+10)}
s_3s_6:
\]
namely,
\[
x=\lim_{n\to\infty}
\widehat{g_1\cdots g_n\xi_1},
\qquad
y=\lim_{n\to\infty}
\widehat{g_1\cdots g_n\xi_3}.
\]
The summable perturbations preserve two extreme directions,
while their first-order terms force every supporting hyperplane
at either limiting point to contain the other point.
\end{remark}

\begin{remark}
All infinite edges in this example have Gram parameter $-2$.
Replacing these entries by $-1$ changes the spectral
decomposition used in the proof: the corresponding rank-two
blocks become affine.
Consequently, the proposition does not answer whether
\[
E_{\mathrm{exp}}(\Phi)=E_{\mathrm{ext}}(\Phi)
\]
holds for every standard Tits realization.
\end{remark}
\section{Facial Structure for Subsets of \texorpdfstring{$E$}{E}}

Let $(\Phi,\Delta)$ be a based root system with associated
Coxeter system $(W,S)$.
Recall that a subset $I\subseteq S$ is called facial if
\[
F_I=\operatorname{conv}(\widehat{\Delta}_I)
\]
is a face of $\operatorname{conv}(\widehat{\Delta})$.

Let $\Phi_I$ denote the corresponding facial root subsystem.

\subsection{Functorial subsets respecting facial structure}

Consider a mapping $E^*$ which associates to every based
root system $(\Phi,\Delta)$ a subset
\[
E^*(\Phi)\subseteq E(\Phi).
\]

We say that $E^*$ is functorial if
\[
E^*(\Phi_I)\subseteq E^*(\Phi)
\]
for every facial root subsystem $\Phi_I$.

We further say that $E^*$ respects the facial structure
if, for every facial subset $I\subseteq S$,
\begin{equation}\label{eq:facial-restriction}
E^*(\Phi_I)=E^*(\Phi)\cap F_I.
\end{equation}

The six subsets constructed in \cite[Section~5]{DHR} satisfy this
property. However, the entire limit-root set $E(\Phi)$
does not satisfy it in general.

The following problem is therefore natural:

\begin{quote}
How can one characterize the functorial subsets of
$E(\Phi)$ which respect the facial structure?
\end{quote}


\par\medskip
\noindent\textbf{Status.} Partially solved. Six compatible constructions are known, but no classification of all such assignments was verified.

\par\medskip
\noindent\textbf{Previous work and verified progress.}
The concrete positive result is \cite[Theorem~5.20]{DHR}, which treats
six natural subsets, including $E_2$. This is not a classification of
all functorial assignments. A claim about one fixed orbit also requires
specifying how its distinguished point is chosen in each subsystem.

\subsection{The single-orbit case}

A first direction would be to investigate the special
case in which
\[
E^*(\Phi)=W\cdot x
\]
for some $x\in E(\Phi)$.


\par\medskip
\noindent\textbf{Status.} No general solution verified. A dense orbit need not equal the limit-root set, and facial compatibility requires an additional argument.

\par\medskip
\noindent\textbf{Previous work and verified progress.}
Minimality gives $\overline{W\cdot x}=E$, not $W\cdot x=E$
\cite[Theorem~3.1]{DHR}. Density therefore does not by itself settle
facial restriction for an orbit. No general single-orbit classification
was located. The assignment must be defined compatibly across root
systems before functoriality can be tested.

\section{Facial Restriction for \texorpdfstring{$E(\Phi)$}{E(Phi)}}

Since $E(\Phi)$ does not respect the facial structure
in the sense of \eqref{eq:facial-restriction}, it is natural to investigate the
relationship between
\[
E(\Phi_I)
\qquad\text{and}\qquad
E(\Phi)\cap F_I
\]
for a facial subset $I\subseteq S$.

\subsection{Joins of orthogonal parabolic limit-root sets}

Suppose that $(\Phi,\Delta)$ is irreducible and that
the facial subsystem $(\Phi_I,\Delta_I)$ is reducible.
Write
\[
\Phi_I=\Phi_1\sqcup\Phi_2,
\]
where
\[
B(\Phi_1,\Phi_2)=0.
\]
Then
\[
E(\Phi_I)=E(\Phi_1)\sqcup E(\Phi_2).
\]

Calculations suggest that, for any
\[
x\in E(\Phi_1),
\qquad
y\in E(\Phi_2),
\]
the line segment joining $x$ and $y$ may be contained
in $E(\Phi)$; that is,
\[
\boxed{
[x,y]\subseteq E(\Phi).
}
\]

This would provide many counterexamples to the facial
restriction formula \eqref{eq:facial-restriction}, provided that both
$E(\Phi_1)$ and $E(\Phi_2)$ are nonempty.


\par\medskip
\noindent\textbf{Status.} \(\checkmark\) Answered in this manuscript. The unrestricted inclusion is false; the theorem below proves it unless both subsystems are non-affine of rank two. This status depends on verification of the proof and counterexample.

\par\medskip
\noindent\textbf{Previous work and verified progress.}
Fu--Reeves--Xu prove that the convex hull of limit roots of mutually
orthogonal irreducible affine standard parabolic subgroups lies in $E$
\cite[Proposition~5.30]{FRX2026}. This includes multiple affine components.
The mixed affine/non-affine cases and the exceptional non-affine dihedral
pair in the theorem below go beyond that cited statement; they are
presented here as manuscript results.

\begin{theorem}\label{thm:orthogonal-parabolic-join}
Assume that $(\Phi,\Pi)$ is irreducible. Let $I,J\subseteq S$
determine irreducible infinite standard parabolic subsystems
satisfying
\[
B(V_I,V_J)=0.
\]
If $\Phi_I$ and $\Phi_J$ are not both non-affine root systems of
rank two, then
\[
E_I*E_J\subseteq E.
\]

For subsets $A,B\subseteq\{v:h(v)=1\}$, define their join by
\[
A*B
=
\{(1-t)x+ty:x\in A,\ y\in B,\ 0\le t\le1\}.
\]

The excluded case cannot be omitted in general: there exists an
irreducible rank-five based root system with two orthogonal
non-affine rank-two standard parabolic subsystems for which
\[
E_I*E_J\not\subseteq E.
\]
\end{theorem}

\begin{remark}\label{rem:meaning-of-exception}
The exception concerns two simultaneously non-affine rank-two
root systems. It does not concern the property of being parabolic:
both subsystems in the theorem are standard parabolic subsystems.

Moreover, the theorem does not assert that the join inclusion
always fails in the excluded case. It asserts that the inclusion
holds outside that case and can fail within it.
\end{remark}

We use two results of Dyer--Hohlweg--Ripoll. First, if a based root
system is irreducible, every orbit in its limit-root set is dense
\cite[Theorem~3.1]{DHR}. Second, an irreducible infinite non-affine
based root system contains a based root subsystem whose simple
roots are positive roots, span the same ambient space, and have
pairwise inner products strictly less than $-1$
\cite[Proposition~6.4]{DHR}.

\begin{lemma}\label{lem:connecting-root}
Under the hypotheses of Theorem~\ref{thm:orthogonal-parabolic-join},
there exists $\rho\in\Phi^+$ such that
\[
\operatorname{supp}(\rho)\subseteq S\setminus(I\cup J),
\]
\[
B(\rho,\alpha)\le0
\qquad
(\alpha\in\Phi_I^+\cup\Phi_J^+),
\]
and neither $B(\rho,-)|_{V_I}$ nor $B(\rho,-)|_{V_J}$ is zero.
Moreover,
\[
w\rho\in\Phi^+
\qquad (w\in W_IW_J).
\]
\end{lemma}

\begin{proof}
Choose a shortest path in the Coxeter graph joining a vertex of
$I$ to a vertex of $J$. Since $B(V_I,V_J)=0$, this path has at least
one internal vertex. By minimality, all its internal vertices lie
outside $I\cup J$.

Starting with the simple root at the first internal vertex,
successively reflect at the remaining internal vertices along
the path. At each step the new simple root has strictly negative
inner product with the root already constructed, so the resulting
root is positive and acquires a strictly positive coefficient at
the new vertex. We obtain a positive root $\rho$ supported
precisely on the internal vertices.

All off-diagonal entries of the Gram matrix are nonpositive.
Consequently,
\[
B(\rho,\alpha)\le0
\qquad
(\alpha\in\Phi_I^+\cup\Phi_J^+).
\]
The first and last edges of the chosen path show that the
restrictions of $B(\rho,-)$ to $V_I$ and $V_J$ are nonzero.

Every element of $W_IW_J$ preserves the coordinates outside
$I\cup J$. At least one such coordinate of $\rho$ is strictly
positive. Hence $w\rho$ cannot be a negative root. The sign
property of roots therefore gives $w\rho\in\Phi^+$.
\end{proof}

\begin{lemma}\label{lem:one-segment-suffices}
Suppose that $x\in E_I$ and $y\in E_J$ satisfy
\[
[x,y]\subseteq E.
\]
Then
\[
E_I*E_J\subseteq E.
\]
\end{lemma}

\begin{proof}
The group $W_I$ fixes $V_J$ pointwise, and $W_J$ fixes $V_I$
pointwise. Thus, for $u\in W_I$ and $v\in W_J$,
\[
(uv)\cdot[x,y]=[u\cdot x,v\cdot y]\subseteq E.
\]
Here the normalized linear action maps the segment onto the
segment between the normalized images of its endpoints.

By minimality,
\[
\overline{W_I\cdot x}=E_I,
\qquad
\overline{W_J\cdot y}=E_J.
\]
Given $x'\in E_I$ and $y'\in E_J$, choose sequences
$u_n\in W_I$ and $v_n\in W_J$ such that
\[
u_n\cdot x\longrightarrow x',
\qquad
v_n\cdot y\longrightarrow y'.
\]
For every $t\in[0,1]$,
\[
(1-t)(u_n\cdot x)+t(v_n\cdot y)\in E.
\]
The closedness of $E$ gives
\[
(1-t)x'+ty'\in E.
\]
\end{proof}

\begin{proof}[Proof of the inclusion in
Theorem~\ref{thm:orthogonal-parabolic-join}]
Fix the positive root $\rho$ supplied by
Lemma~\ref{lem:connecting-root}. By
Lemma~\ref{lem:one-segment-suffices}, it is enough to construct one
segment $[x,y]\subseteq E$ with $x\in E_I$ and $y\in E_J$.

\medskip
\noindent
\textit{Case 1: at least one subsystem is affine.}

Suppose first that $\Phi_I$ is affine. Let
\[
E_I=\{\delta\},
\qquad h(\delta)=1.
\]
The affine root-string structure provides a root $a$ and a
constant $\ell>0$ such that
\[
a_n=a+n\ell\delta\in\Phi_I^+
\qquad(n\ge0).
\]
Since $\delta$ has strictly positive coordinates on $I$,
Lemma~\ref{lem:connecting-root} implies
\[
B(\rho,\delta)<0.
\]
Consequently,
\[
\begin{aligned}
s_{a_n}\rho
&=\rho-2B(\rho,a_n)a_n\\
&=Cn^2\delta+O(n),
\end{aligned}
\]
where
\[
C=-2\ell^2B(\rho,\delta)>0.
\]

Assume that $\Phi_J$ is non-affine. By
\cite[Proposition~6.4]{DHR}, we can choose positive roots
$c,d\in\Phi_J^+$ such that
\[
B(c,d)<-1,
\qquad B(\rho,c)<0.
\]
Let $t=s_cs_d$, and let $y\in E_J$ be its normalized attracting
eigenvector, with eigenvalue $\mu>1$. Since
\[
B(\rho,c)<0,\qquad B(\rho,d)\le0,
\]
the attracting coefficient of $\rho$ is strictly positive.
Thus, for some $D>0$,
\[
t^m\rho=D\mu^m y+O(1).
\]
The element $t$ fixes $V_I$ pointwise, so
\[
s_{a_n}t^m\rho
=
Cn^2\delta+D\mu^m y+O(n+1).
\]
All these vectors are positive roots by
Lemma~\ref{lem:connecting-root}.

Given $r>0$, choose positive integers $n_m$ satisfying
\[
n_m\sim\sqrt{\frac{rD}{C}\mu^m}.
\]
Then
\[
\widehat{s_{a_{n_m}}t^m\rho}
\longrightarrow
\frac{r\delta+y}{r+1}.
\]
Hence $[\delta,y]\subseteq E$.

If $\Phi_J$ is also affine, write
\[
E_J=\{\varepsilon\},
\qquad h(\varepsilon)=1,
\]
and choose an affine root string
\[
b_m=b+m\ell'\varepsilon\in\Phi_J^+.
\]
Orthogonality gives
\[
s_{a_n}s_{b_m}\rho
=
Cn^2\delta+Dm^2\varepsilon+O(n+m+1)
\]
for constants $C,D>0$.
For any $r>0$, choosing $n,m\to\infty$ with
\[
\frac{Cn^2}{Dm^2}\longrightarrow r
\]
yields
\[
\widehat{s_{a_n}s_{b_m}\rho}
\longrightarrow
\frac{r\delta+\varepsilon}{r+1}.
\]
Thus $[\delta,\varepsilon]\subseteq E$, completing this case.

\medskip
\noindent
\textit{Case 2: both subsystems are non-affine.}

By hypothesis, one subsystem has rank at least three.
Interchanging $I$ and $J$ if necessary, assume
\[
\dim V_I\ge3.
\]
Using \cite[Proposition~6.4]{DHR}, choose linearly independent
positive roots $a,b,u\in\Phi_I^+$ that belong to the simple system
of a based root subsystem and satisfy
\[
B(a,b),B(a,u),B(b,u)<-1,
\qquad B(\rho,a)<0.
\]
Likewise choose $c,d\in\Phi_J^+$ with
\[
B(c,d)<-1,\qquad B(\rho,c)<0.
\]

Set
\[
g=s_as_b,\qquad t=s_cs_d.
\]
Choose normalized eigenvectors
\[
gx=\lambda x,\qquad gx_-=\lambda^{-1}x_-,
\]
\[
ty=\mu y,\qquad ty_-=\mu^{-1}y_-,
\]
where
\[
\lambda,\mu>1,\qquad
x,x_-\in E_I,\qquad y,y_-\in E_J.
\]
The planes
\[
L=\operatorname{span}\{a,b\},
\qquad
M=\operatorname{span}\{c,d\}
\]
are nondegenerate and mutually orthogonal. Every $v\in V$
therefore has a unique decomposition
\[
v=A(v)x+A_-(v)x_-+D(v)y+D_-(v)y_-+v_0,
\]
where
\[
v_0\in(L\oplus M)^\perp.
\]
The connecting-root inequalities imply
\[
A(\rho),A_-(\rho),D(\rho),D_-(\rho)>0.
\]

Suppose first that
\[
\frac{\log\lambda}{\log\mu}\notin\mathbb Q.
\]
For every $r>0$, irrational rotation on the circle gives sequences
$n_k,m_k\to\infty$ such that
\[
\frac{A(\rho)\lambda^{n_k}}
     {D(\rho)\mu^{m_k}}
\longrightarrow r.
\]
Since
\[
g^{n_k}t^{m_k}\rho
=
A(\rho)\lambda^{n_k}x
+D(\rho)\mu^{m_k}y+O(1),
\]
we obtain
\[
\widehat{g^{n_k}t^{m_k}\rho}
\longrightarrow
\frac{rx+y}{r+1}.
\]
Thus $[x,y]\subseteq E$.

It remains to consider
\[
\frac{\log\lambda}{\log\mu}\in\mathbb Q.
\]
Replacing $g$ and $t$ by suitable positive powers, we may assume
that their attracting eigenvalues coincide:
\[
\lambda=\mu=\Lambda>1.
\]

Define
\[
\mathcal R
=
\left\{
r>0:\frac{rx+y}{r+1}\in E
\right\}.
\]
This is a closed subset of $(0,\infty)$. Since $g$ fixes $y$ and
sends $x$ to $\Lambda x$, the $W$-invariance of $E$ gives
\[
\Lambda^{\mathbb Z}\mathcal R=\mathcal R.
\]

Put
\[
\kappa=-B(x,x_-)>0.
\]
The inequalities $B(u,a),B(u,b)<-1$ imply
\[
A(u)>0,\qquad A_-(u)>0.
\]
Set
\[
d_0=\kappa A(u)A_-(u)>0,
\qquad
u_m=g^m u.
\]
Because $a,b,u$ belong to a simple system of a based root
subsystem, $u_m\in\Phi_I^+$ for $m\ge0$.

Writing
\[
u=A(u)x+A_-(u)x_-+u_0,
\qquad u_0\perp L,
\]
and using $B(u,u)=1$, we obtain
\[
B(u,u_m)
=
1+2d_0-d_0(\Lambda^m+\Lambda^{-m}).
\]
For all sufficiently large $m$, therefore,
\[
q_m:=-B(u,u_m)>1.
\]
Let
\[
H_m=s_{u_m}s_u.
\]
Its attracting eigenvalue is
\[
\eta_m=
\left(q_m+\sqrt{q_m^2-1}\right)^2.
\]

For sufficiently large $m$,
\[
B(\rho,u)\le0,\qquad B(\rho,u_m)<0.
\]
Both isotropic eigenvectors of $H_m$ in
$\operatorname{span}\{u,u_m\}$ are positive linear combinations
of $u$ and $u_m$. It follows that the attracting coefficient of
$\rho$ is strictly positive.

Moreover,
\[
A(u)>0,\qquad A(u_m)=\Lambda^mA(u)>0.
\]
Applying the linear functional $A$ to the eigenspace decomposition
of $H_m^k\rho$ consequently gives
\[
A(H_m^k\rho)=C_m\eta_m^k+O(1)
\qquad(k\to\infty)
\]
for some $C_m>0$.

Since $H_m\in W_I$,
\[
D(H_m^k\rho)=D(\rho)=:D_0>0.
\]
Fixing a sufficiently large $k$ and letting $N\to\infty$, we have
\[
\widehat{g^Nt^NH_m^k\rho}
\longrightarrow
\frac{A(H_m^k\rho)x+D_0y}
     {A(H_m^k\rho)+D_0}.
\]
Hence
\[
\frac{A(H_m^k\rho)}{D_0}\in\mathcal R.
\]

Consider the closed subset
\[
\mathcal T
=
\{\log_\Lambda r+\mathbb Z:r\in\mathcal R\}
\subseteq\mathbb R/\mathbb Z.
\]
Its closedness follows from the closedness of $\mathcal R$ and
the invariance $\Lambda^{\mathbb Z}\mathcal R=\mathcal R$.
For every sufficiently large $m$, the set $\mathcal T$ contains
all accumulation points of
\[
k\log_\Lambda\eta_m
+\log_\Lambda(C_m/D_0)+o(1)
\pmod{\mathbb Z}.
\]

If some $\log_\Lambda\eta_m$ is irrational, these accumulation
points fill the circle, and hence
\[
\mathcal T=\mathbb R/\mathbb Z.
\]
If all $\log_\Lambda\eta_m$ are rational but their reduced
denominators are unbounded, then $\mathcal T$ contains translates
of equally spaced grids of arbitrarily large cardinality.
Closedness again gives
\[
\mathcal T=\mathbb R/\mathbb Z.
\]

Suppose, finally, that all these numbers are rational with
uniformly bounded reduced denominators. Then there exists
$M\in\mathbb N$ such that
\[
\log_\Lambda\eta_m-2m\in\frac1M\mathbb Z
\]
for every sufficiently large $m$.

On the other hand,
\[
q_m=d_0\Lambda^m-(1+2d_0)+d_0\Lambda^{-m},
\]
so
\[
\frac{\eta_m}{\Lambda^{2m}}\longrightarrow4d_0^2.
\]
Furthermore,
\[
q_m<d_0\Lambda^m,
\qquad
\eta_m<4q_m^2,
\]
and therefore
\[
\frac{\eta_m}{\Lambda^{2m}}<4d_0^2.
\]
Thus the sequence
\[
\log_\Lambda\eta_m-2m
\]
lies in the discrete set $\frac1M\mathbb Z$, converges to
$\log_\Lambda(4d_0^2)$, and is strictly smaller than its limit.
This is impossible.

Consequently,
\[
\mathcal T=\mathbb R/\mathbb Z.
\]
Together with $\Lambda^{\mathbb Z}\mathcal R=\mathcal R$, this
implies
\[
\mathcal R=(0,\infty).
\]
Hence $[x,y]\subseteq E$. Applying
Lemma~\ref{lem:one-segment-suffices} completes the proof of the
join inclusion.
\end{proof}

\begin{proof}[Construction and verification of the counterexample]
Let $\Pi=\{\alpha_1,\ldots,\alpha_5\}$, and take the Gram matrix
\[
G=
\begin{pmatrix}
1&-2&0&0&-2\\
-2&1&0&0&-2\\
0&0&1&-2&-2\\
0&0&-2&1&-2\\
-2&-2&-2&-2&1
\end{pmatrix}.
\]
The Coxeter graph is connected, so the ambient root system is
irreducible. Set
\[
I=\{1,2\},\qquad J=\{3,4\}.
\]
Then $V_I\perp V_J$, and both standard parabolic subsystems are
non-affine of rank two.

Put
\[
\tau=2-\sqrt3,
\qquad
x=\frac{\alpha_1+\tau\alpha_2}{1+\tau},
\qquad
y=\frac{\alpha_3+\tau\alpha_4}{1+\tau}.
\]
The matrix of $s_1s_2$ on $V_I$ is
\[
\begin{pmatrix}
15&-4\\
4&-1
\end{pmatrix},
\]
whose attracting eigenvalue is $\tau^{-2}$ and whose attracting
eigenvector is $(1,\tau)$. Hence
\[
\widehat{(s_1s_2)^n\alpha_1}\longrightarrow x.
\]
Similarly,
\[
\widehat{(s_3s_4)^n\alpha_3}\longrightarrow y.
\]
Thus
\[
x\in E_I,\qquad y\in E_J.
\]

We shall prove that
\[
z=\frac{x+2y}{3}\notin E.
\]

Define
\[
p=\frac{1+\sqrt3}{2},
\qquad
q=\frac{\sqrt3-1}{2}.
\]
For
\[
v=c_1\alpha_1+c_2\alpha_2+c_3\alpha_3+c_4\alpha_4+d\alpha_5,
\]
set
\[
\begin{aligned}
A_+(v)&=pc_1-qc_2+2d,
&
A_-(v)&=-qc_1+pc_2+2d,\\
C_+(v)&=pc_3-qc_4+2d,
&
C_-(v)&=-qc_3+pc_4+2d.
\end{aligned}
\]
Consider the homogeneous quadratic polynomial
\[
H(v)
=
\bigl(20A_+(v)-9C_+(v)\bigr)
\bigl(3A_+(v)-2C_+(v)\bigr).
\]
We claim that
\[
H(\beta)\ge0
\qquad(\beta\in\Phi^+).
\]

All roots have integral coordinates in the basis $\Pi$.
If $\beta=\sum_i b_i\alpha_i$ is a positive root that is not
simple, then
\[
1=B(\beta,\beta)=\sum_i b_iB(\beta,\alpha_i)
\]
provides an index $i$ with
\[
b_i>0,\qquad B(\beta,\alpha_i)>0.
\]
The root $s_i\beta$ remains positive: another positive coordinate
of $\beta$ is unchanged, so $s_i\beta$ cannot be negative.
Moreover,
\[
s_i\beta=\beta-2B(\beta,\alpha_i)\alpha_i
\]
has strictly smaller height and differs from $\beta$ only in its
$i$-th coordinate.

Iterating this descent reaches a simple root, since the heights
are positive integers. Reversing the descent shows that every
positive root admits a construction from a simple root along
which all coordinates are nondecreasing.

Fix such a construction of $\beta$, and suppose first that its
fifth coordinate $d$ is positive. If the construction uses $s_5$,
let $\rho$ be the root immediately after the last application of
$s_5$. Write
\[
\rho=a_1\alpha_1+\cdots+a_4\alpha_4+d\alpha_5,
\qquad
S=a_1+\cdots+a_4.
\]
If $d_{\mathrm{old}}$ is the fifth coordinate immediately before
that reflection, then
\[
d=4S-d_{\mathrm{old}}.
\]
Coordinate monotonicity gives $d\ge d_{\mathrm{old}}$, and hence
\[
d\ge2S.
\]
If the construction never uses $s_5$, its initial simple root
must be $\alpha_5$; in this case take $\rho=\alpha_5$, and the same
inequality holds.

In either case,
\[
\beta=w\rho
\qquad\text{for some }w\in W_IW_J,
\]
and
\[
a_i\ge0,\qquad \sum_{i=1}^4a_i\le\frac d2.
\]
Consequently,
\[
Ld\le A_\pm(\rho),C_\pm(\rho)\le Ud,
\]
where
\[
L=2-\frac q2=\frac{9-\sqrt3}{4},
\qquad
U=2+\frac p2=\frac{9+\sqrt3}{4}.
\]
Set
\[
R=\frac UL=\frac{9+\sqrt3}{9-\sqrt3}<\frac32.
\]

A direct calculation gives
\[
\begin{aligned}
(A_+,A_-)(s_1v)
&=(\tau^{-1}A_-(v),\,\tau A_+(v)),\\
(A_+,A_-)(s_2v)
&=(\tau A_-(v),\,\tau^{-1}A_+(v)).
\end{aligned}
\]
Both $s_1$ and $s_2$ fix $C_+$ and $C_-$. Likewise,
\[
\begin{aligned}
(C_+,C_-)(s_3v)
&=(\tau^{-1}C_-(v),\,\tau C_+(v)),\\
(C_+,C_-)(s_4v)
&=(\tau C_-(v),\,\tau^{-1}C_+(v)),
\end{aligned}
\]
and $s_3,s_4$ fix $A_+,A_-$.

It follows that
\[
A_+(\beta),C_+(\beta)>0,
\qquad
\frac{A_+(\beta)}{C_+(\beta)}
\in
\bigcup_{m\in\mathbb Z}\tau^m[R^{-1},R].
\]
Since $\tau<3/10$ and $R<3/2$, we have
\[
\tau R<\frac9{20},
\qquad
R^{-1}>\frac23.
\]
For $m\ge1$,
\[
\tau^mR\le\tau R<\frac9{20},
\]
whereas for $m\le0$,
\[
\tau^mR^{-1}\ge R^{-1}>\frac23.
\]
Thus the displayed union avoids the interval
$(9/20,2/3)$. Therefore
\[
\left(
20\frac{A_+(\beta)}{C_+(\beta)}-9
\right)
\left(
3\frac{A_+(\beta)}{C_+(\beta)}-2
\right)
\ge0,
\]
which proves $H(\beta)\ge0$ when $d>0$.

Now suppose that the fifth coordinate of $\beta$ is zero.
Its coordinatewise nondecreasing construction cannot use $s_5$.
Since $W_I$ and $W_J$ preserve their respective summands, the
root $\beta$ belongs to $\Phi_I$ or $\Phi_J$. In the first case
$C_+(\beta)=0$, and in the second case $A_+(\beta)=0$. Accordingly,
\[
H(\beta)=60A_+(\beta)^2
\qquad\text{or}\qquad
H(\beta)=18C_+(\beta)^2.
\]
This proves the claim for every positive root.

The homogeneity of $H$ implies
\[
H(\widehat\beta)=\frac{H(\beta)}{h(\beta)^2}\ge0.
\]
Taking accumulation points and using continuity, we obtain
\[
E\subseteq\{v\in V:H(v)\ge0\}.
\]

Finally, direct substitution gives
\[
A_+(x)=1,\qquad C_+(x)=0,
\qquad
A_+(y)=0,\qquad C_+(y)=1.
\]
Hence
\[
A_+(z)=\frac13,\qquad C_+(z)=\frac23,
\]
and
\[
H(z)
=
\left(\frac{20}{3}-6\right)
\left(1-\frac43\right)
=
-\frac29<0.
\]
Therefore $z\notin E$, although $z\in[x,y]$ with
$x\in E_I$ and $y\in E_J$. This proves
\[
E_I*E_J\not\subseteq E.
\]
\end{proof}

\subsection{Join inclusion in the standard Tits realization}

Does $E_I*E_J\subseteq E$ always hold for orthogonal irreducible
infinite standard parabolic subsystems in the standard Tits realization?


\par\medskip
\noindent\textbf{Status.} \(\checkmark\) Answered in this manuscript, conditional on its join theorem. Infinite rank-two standard parabolic subsystems in the standard Tits realization are affine, excluding the exceptional configuration.

\par\medskip
\noindent\textbf{Previous work and verified progress.}
The affine-component case is covered by
\cite[Proposition~5.30]{FRX2026}. This does not establish the full statement
below when a standard parabolic factor is non-affine of higher rank.
The corollary below instead uses this manuscript's join theorem and
the fact that infinite rank-two subsystems in the standard Tits
realization are affine.

\begin{corollary}\label{cor:standard-tits-join}
Suppose that $(\Phi,\Pi)$ is an irreducible root system in the
standard Tits realization. If $\Phi_I,\Phi_J$ are orthogonal
irreducible infinite standard parabolic subsystems, then
\[
E_I*E_J\subseteq E.
\]
\end{corollary}

\begin{proof}
In the standard Tits realization, an infinite edge has inner
product $-1$. Every infinite rank-two standard parabolic
subsystem is therefore affine. The exceptional case in
Theorem~\ref{thm:orthogonal-parabolic-join} cannot occur.
\end{proof}

\begin{remark}\label{rem:reducible-parabolic-join}
For reducible standard parabolic subsystems, the theorem must be
applied to pairs of irreducible components.

More precisely, suppose
\[
\Phi_I=\bigsqcup_a\Phi_{I_a},
\qquad
\Phi_J=\bigsqcup_b\Phi_{J_b}.
\]
Then
\[
E_I=\bigcup_a E_{I_a},
\qquad
E_J=\bigcup_b E_{J_b}.
\]
If no pair of infinite components
$(\Phi_{I_a},\Phi_{J_b})$ consists of two non-affine rank-two
root systems, then
\[
E_I*E_J
=
\bigcup_{a,b}(E_{I_a}*E_{J_b})
\subseteq E.
\]
In particular, having total rank greater than two does not by
itself exclude the exceptional configuration in a reducible
subsystem.
\end{remark}

\begin{remark}\label{rem:convex-hull-isotropic-counterexample}
In the rank-five counterexample,
\[
B(x,x)=B(y,y)=B(x,y)=0.
\]
Thus the entire segment $[x,y]$ lies in the isotropic cone
\[
Q=\{v\in V:B(v,v)=0\}.
\]
The point $z=(x+2y)/3$ therefore satisfies
\[
z\in\operatorname{conv}(E)\cap Q,
\qquad z\notin E.
\]
Consequently,
\[
\operatorname{conv}(E)\cap Q\ne E
\]
for this irreducible based root system.

This construction does not assert that $x$ and $y$ are
accumulation points of a common dominance chain.
\end{remark}

\subsection{The proposed join formula for facial intersections}

Let $\Phi_I=\bigsqcup_{k=1}^{m}\Phi_{I,k}$ be the decomposition
into irreducible components, and put $V_I=\operatorname{span}(\Delta_I)$.
Is it always true that
\[
E(\Phi)\cap V_I=\operatorname{join}_{k=1}^{m}E(\Phi_{I,k})?
\]
The counterexample in Proposition~\ref{prop:counterexample-convE-Q}
already gives a negative answer. Take $I=\{1,2,3,4\}$.
The point $z=(x+2y)/3$ belongs to the join of the two component
limit-root sets, but does not belong to $E(\Phi)\cap V_I$.


\par\medskip
\noindent\textbf{Status.} \(\checkmark\) Answered negatively in this manuscript. Its counterexample disproves the proposed join formula; this is related to the same construction used in Question~1.1.

\par\medskip
\noindent\textbf{Previous work and verified progress.}
Hohlweg--Labb\'e--Ripoll give a rank-five example in which
$E(\Phi_I)\subsetneq E(\Phi)\cap F_I$
\cite[Example~5.8]{HLR}. That example disproves ordinary facial
restriction, not the stronger join formula. The negative answer to the
join formula given above uses this manuscript's different counterexample.
These two failures should not be conflated.

\subsection{Obstructions to facial restriction}
How can $E(\Phi)\cap V_I$ be described exactly? Is mixing between
irreducible components the only obstruction to
\eqref{eq:facial-restriction} for the full limit-root set?
This question remains open in this manuscript.


\par\medskip
\noindent\textbf{Status.} No general solution verified. Existing obstructions do not classify every failure of facial restriction.

\par\medskip
\noindent\textbf{Previous work and verified progress.}
For the affine part, \cite[Corollary~5.33]{FRX2026} identifies
$E\cap Z$ through the Pos-finite condition, where $Z$ is the normalized
\emph{unclosed} imaginary cone. This is useful for separating affine
combinations from other limit roots, but it does not describe the entire
intersection $E\cap V_I$. No proof that component mixing is the only
obstruction was located.

\subsection{Facial restriction for irreducible subsystems}

If $\Phi_I$ is irreducible, is it always true that
\[
E(\Phi_I)=E(\Phi)\cap F_I?
\]

\par\medskip
\noindent\textbf{Status.}
\(\checkmark\) \textit{Answered negatively in this manuscript.}
The rank-six counterexample below has an irreducible rank-five
facial subsystem for which the equality fails, even in the
relative interior of the face.

\par\medskip
\noindent\textbf{Previous work and verified progress.}
This question is raised in \cite[Section~7.3]{DHR}.
The following construction extends
Proposition~\ref{prop:counterexample-convE-Q}.
An additional simple root produces a sequence of roots converging
to a point excluded from the limit-root set of the original
irreducible subsystem.

\begin{proposition}\label{prop:irreducible-facial-counterexample}
There exists an irreducible based root system $\Phi$ of rank six
and signature $(4,2)$ with an irreducible standard parabolic
subsystem $\Phi_I$ of rank five and signature $(3,2)$ such that
\[
E(\Phi_I)\subsetneq E(\Phi)\cap F_I.
\]
Moreover,
\[
\bigl(E(\Phi)\cap\operatorname{relint}(F_I)\bigr)
\setminus E(\Phi_I)\ne\varnothing.
\]
\end{proposition}

\begin{proof}
Let $V=\bigoplus_{i=1}^{6}\mathbb R\alpha_i$, and define $B$
by the Gram matrix
\[
G=
\begin{pmatrix}
1&-2&0&0&-2&-1\\
-2&1&0&0&-2&-1\\
0&0&1&-2&-2&-2\\
0&0&-2&1&-2&-2\\
-2&-2&-2&-2&1&0\\
-1&-1&-2&-2&0&1
\end{pmatrix}.
\]
Let
\[
s_i(v)=v-2B(v,\alpha_i)\alpha_i,
\qquad
W=\langle s_1,\ldots,s_6\rangle,
\qquad
\Phi=W\{\alpha_1,\ldots,\alpha_6\}.
\]
All off-diagonal entries of $G$ are zero or at most $-1$,
so these data define a based root system.

Put
\[
I=\{1,2,3,4,5\},
\qquad
h\left(\sum_{i=1}^{6}c_i\alpha_i\right)=\sum_{i=1}^{6}c_i,
\qquad
\widehat v=\frac{v}{h(v)}.
\]
Since the simple roots form a basis,
\[
F_I=\operatorname{conv}\{\alpha_1,\ldots,\alpha_5\}
\]
is a face of the simple-root simplex.
The fifth vertex is joined to each of the first four vertices,
so $\Phi_I$ is irreducible. The sixth vertex is also joined to
the first four vertices, so $\Phi$ is irreducible.

The restriction of $G$ to $I$ is precisely the Gram matrix in
Proposition~\ref{prop:counterexample-convE-Q}.
Set
\[
\tau=2-\sqrt3,
\qquad
\lambda=\tau^{-2}=7+4\sqrt3,
\]
and define
\[
x_+=\frac{\alpha_1+\tau\alpha_2}{1+\tau},
\qquad
x_-=\frac{\tau\alpha_1+\alpha_2}{1+\tau},
\]
\[
y_+=\frac{\alpha_3+\tau\alpha_4}{1+\tau},
\qquad
y_-=\frac{\tau\alpha_3+\alpha_4}{1+\tau}.
\]
The exclusion established in that proposition gives
\begin{equation}\label{eq:facial-excluded-point}
z:=\frac{x_++2y_+}{3}\notin E(\Phi_I).
\end{equation}

We show that $z\in E(\Phi)$.
Let
\[
w=s_1s_2s_3s_4,
\qquad
u=\alpha_6-\alpha_1-\alpha_2-2\alpha_3-2\alpha_4.
\]
The blocks $\{1,2\}$ and $\{3,4\}$ are orthogonal, and direct
calculation gives
\[
wx_+=\lambda x_+,
\qquad
wx_-=\lambda^{-1}x_-,
\qquad
wy_+=\lambda y_+,
\qquad
wy_-=\lambda^{-1}y_-.
\]
Also,
\[
B(u,\alpha_i)=0
\qquad(1\le i\le4),
\]
so $wu=u$. Since
\[
x_++x_-=\alpha_1+\alpha_2,
\qquad
y_++y_-=\alpha_3+\alpha_4,
\]
we obtain
\[
\alpha_6=u+x_++x_-+2y_++2y_-.
\]
Consequently,
\begin{equation}\label{eq:facial-root-iterate}
w^n\alpha_6
=
u+\lambda^n(x_++2y_+)
+\lambda^{-n}(x_-+2y_-)
\qquad(n\ge0).
\end{equation}

Each $w^n\alpha_6$ is a positive root.
Indeed, reflections indexed by $\{1,2,3,4\}$ leave its sixth
coordinate unchanged, so that coordinate is always $1$.
The root therefore cannot be negative.

As $h(u)=-5$ and $h(x_\pm)=h(y_\pm)=1$,
\[
h(w^n\alpha_6)=3\lambda^n+3\lambda^{-n}-5.
\]
These heights are strictly increasing for $n\ge0$ and tend to
infinity. Thus the roots are distinct, and
\eqref{eq:facial-root-iterate} yields
\[
\widehat{w^n\alpha_6}
=
\frac{
u+\lambda^n(x_++2y_+)+\lambda^{-n}(x_-+2y_-)
}{
3\lambda^n+3\lambda^{-n}-5
}
\longrightarrow z.
\]
Hence
\[
z\in E(\Phi)\cap F_I.
\]
Together with \eqref{eq:facial-excluded-point} and the inclusion
\[
E(\Phi_I)\subseteq E(\Phi)\cap F_I,
\]
this proves strict inclusion.

To obtain a counterexample in the relative interior of $F_I$,
observe that the first four coordinates of $z$ are positive
and sum to $1$, whereas its fifth and sixth coordinates vanish.
Thus
\[
B(z,\alpha_5)=-2,
\qquad
s_5z=z+4\alpha_5,
\qquad
h(s_5z)=5.
\]
Define
\begin{equation}\label{eq:facial-interior-counterexample}
z^\circ:=\widehat{s_5z}
=\frac{z+4\alpha_5}{5}.
\end{equation}
All coordinates of $z^\circ$ indexed by $I$ are strictly positive,
so
\[
z^\circ\in\operatorname{relint}(F_I).
\]
The normalized $W$-action preserves $E(\Phi)$, hence
$z^\circ\in E(\Phi)$.
If $z^\circ\in E(\Phi_I)$, normalized $W_I$-invariance would give
\[
\widehat{s_5z^\circ}=z\in E(\Phi_I),
\]
contradicting \eqref{eq:facial-excluded-point}.
Therefore
\[
z^\circ\in
\bigl(E(\Phi)\cap\operatorname{relint}(F_I)\bigr)
\setminus E(\Phi_I).
\]

Finally, the signatures can be checked explicitly.
Put
\[
\eta=\alpha_5-2\alpha_1-2\alpha_2-2\alpha_3-2\alpha_4.
\]
In the basis $(x_+,x_-,y_+,y_-,\eta,u)$, the Gram matrix is
\[
\begin{pmatrix}
0&-1\\
-1&0
\end{pmatrix}
\oplus
\begin{pmatrix}
0&-1\\
-1&0
\end{pmatrix}
\oplus
\begin{pmatrix}
17&12\\
12&11
\end{pmatrix}.
\]
The last block is positive definite, since its determinant is
$43$ and its first diagonal entry is positive.
Thus $B$ has signature $(4,2)$.
The first five basis vectors span $V_I$, on which the signature
is $(3,2)$.
\end{proof}

\begin{remark}
Although the initial counterexample point $z$ has reducible
support $\{1,2,3,4\}$, the point $z^\circ$ in
\eqref{eq:facial-interior-counterexample} has support exactly $I$.
Thus the equality can fail even when the facial subsystem is
irreducible and the counterexample point has full support in it.

This construction uses nonstandard Gram parameters on infinite
edges. It does not settle the corresponding question restricted
to the standard Tits realization.
\end{remark}

\section{Relation with Hyperbolic Geometry and Geometric Group Theory}

The relationship with hyperbolic geometry is well developed in the
Lorentzian case; extensions to general signatures remain a research direction.

Consider a based root system $(\Phi,\Delta)$ of rank
three or four, of indefinite and weakly hyperbolic type.

In this situation, the Coxeter group $W$ acts on
$E(\Phi)$ as a group generated by hyperbolic reflections.
Consequently, it can be viewed as a Fuchsian group or
a Kleinian group.

Moreover, the limit-root set $E(\Phi)$ coincides with
the limit set $\Lambda(W)$ of $W$ viewed as a Kleinian
group.

This correspondence suggests a natural research
direction: to generalize classical properties of limit
sets of Kleinian groups to the more general setting
of Coxeter groups acting on quadratic spaces.

\subsection{Dynamics of the action}
Which classical dynamical properties of Kleinian limit sets extend
to the action of $W$ on $E(\Phi)$?


\par\medskip
\noindent\textbf{Status.} Research direction with established partial results. Known dynamical theorems do not settle every proposed extension.

\par\medskip
\noindent\textbf{Previous work and verified progress.}
A higher-signature development is the theory of convex-cocompact
projective reflection representations \cite{DGKLM2025}. For word-hyperbolic
Coxeter groups, the representations singled out there are projective
Anosov. These results give boundary dynamics beyond classical hyperbolic
space, but require the paper's conditions on the Cartan matrix.

\subsection{Geometry of the limit set}
Which geometric properties of Kleinian limit sets extend to
limit-root sets in general quadratic spaces?


\par\medskip
\noindent\textbf{Status.} Research direction with established partial results. Lorentzian and convex projective methods cover substantial special cases.

\par\medskip
\noindent\textbf{Previous work and verified progress.}
For spectral and arrangement methods, see \cite[Sections~3--4]{ChenLabbe2017}.
For a projective convex framework, see \cite{DGKLM2025}.
These references do not supply a classification in every signature.


\subsection{The imaginary convex body and the convex core}
What is the relationship between the imaginary convex body and
the convex core in the hyperbolic setting, and how can it be generalized?


\par\medskip
\noindent\textbf{Status.} Solved in the Lorentzian setting; partially solved in general. Extensions of the convex-core interpretation to other signatures require additional arguments.

\par\medskip
\noindent\textbf{Previous work and verified progress.}
Dyer--Hohlweg--Ripoll explain the Lorentzian convex-core interpretation
in \cite[Section~7.4]{DHR}. For projective reflection groups,
\cite[Section~5]{DGKLM2025} studies the minimal invariant convex subset
of the Tits--Vinberg domain. Keep the imaginary cone, its closure,
and the convex core quotient distinct when comparing the constructions.

\subsection{Reflection groups of infinite covolume}
How do these relationships extend to discrete reflection groups
of infinite covolume?


\par\medskip
\noindent\textbf{Status.} Lorentzian limit-set identification solved. The equality between root and hyperbolic limit sets includes infinite-covolume groups. Further geometric conclusions require separate hypotheses.

\par\medskip
\noindent\textbf{Previous work and verified progress.}
The Lorentzian equality between root and hyperbolic limit sets
\cite{HPR2020} is not confined to lattices. The Cartan-matrix criteria
of \cite[Theorems~1.3 and~1.8]{DGKLM2025} provide further reflection-group
examples and a projective replacement for classical convex cocompactness.
They do not assert convex cocompactness for every infinite-covolume group.

\section{Dynamics of the Projective Action of \texorpdfstring{$W$}{W}}

\subsection{Orbit limit sets of arbitrary directions}

Another natural question concerns the dynamics of the
projective action of $W$ on all directions of the vector
space $V$, rather than only on the directions of roots
and the imaginary cone.

For weakly hyperbolic root systems, previous work of
Chen and Labbé shows that $E(\Phi)$ coincides with
the set of limit directions arising from the projective
action on the weights of the root system.

However, limit sets of orbits of other directions may
also contain points outside $E(\Phi)$.

Thus, it is natural to investigate the general structure
of the limit sets arising from the projective action
of $W$ on arbitrary directions of $V$.


\par\medskip
\noindent\textbf{Status.} Partially solved. Lorentzian orbit accumulation has substantial results, but no complete classification for arbitrary directions and signatures was verified.

\par\medskip
\noindent\textbf{Previous work and verified progress.}
Chen--Labb\'e prove that any light-like accumulation direction of
an injective orbit sequence in a Lorentzian system is a limit root
\cite[Theorem~2.6]{ChenLabbe2017}. Their Section~4 also bounds the full set of
limit directions using unimodular eigenspaces and the reflection
arrangement; Section~4.4 retains further open questions even there.
Thus this subsection has substantial partial answers, not a complete
classification of all orbit closures in arbitrary signature.

\section{Ergodic Theory for the \texorpdfstring{$W$}{W}-Action on \texorpdfstring{$E$}{E}}

A classical problem in the ergodic theory of discrete
groups is to determine whether a group acting on its
limit set admits an invariant measure, whether such
a measure is unique, and how to construct it.

In the present setting, a natural problem is to study
the existence and properties of $W$-invariant measures
on $E(\Phi)$.

\begin{remark}
In \cite[Section~7.6]{DHR}, the authors ask about invariant measures
on $E$ without specifying their total mass. In the next three
subsections we classify $W$-invariant Borel probability measures.
Since $E$ is compact and metrizable, normalization identifies these
with nonzero invariant Radon measures up to positive scalar multiples.
If infinite mass is allowed, every orbit $Wx$ carries the invariant
counting measure $A\mapsto\#(A\cap Wx)$, whose total mass is infinite
when $Wx$ is infinite.
\end{remark}

\subsection{Existence of invariant measures}
Throughout these three subsections, $\Phi$ is a finite-rank based root
system and $W$ acts on $E=E(\Phi)$ by normalization.
A Borel probability measure $\mu$ is $W$-invariant if
$\mu(w\cdot A)=\mu(A)$ for every Borel set $A\subseteq E$ and every
$w\in W$.

\noindent\textbf{Status.} $\checkmark$ Solved. An invariant probability
exists precisely when $E$ contains a finite $W$-orbit.

\begin{proposition}\label{prop:limit-root-invariant-probability}
If $\Phi$ is irreducible of indefinite type and has rank at least
three, then $E$ admits no $W$-invariant Borel probability measure.
\end{proposition}

\begin{proof}
Suppose that $\mu$ is such a measure. By
\cite[Theorem~3.1(b)]{DHR}, every orbit is dense in $E$. In particular,
the support of $\mu$, being nonempty, closed and invariant, equals $E$.
The infinite group $W$ acts faithfully on $E$ by
\cite[Theorem~6.1]{DHR}, so $E$ is infinite. A finite orbit would be
closed and dense, hence every orbit is infinite. Consequently
$\mu(\{x\})=0$ for each $x\in E$, since all points of its orbit have
the same mass.

Fix $x\in E$. By \cite[Lemma~3.8(b)]{DHR}, there is $y\in E$ with
$B(x,y)<0$. Choose $0<\varepsilon<-B(x,y)$ and set
\[
K=\{z\in E:B(x,z)\le -\varepsilon\}.
\]
This compact set contains a relative neighborhood of $y$, so
$\mu(K)>0$. Choose positive roots $\alpha_n$ with
$\widehat{\alpha_n}\to x$. Put $Z=\operatorname{conv}(E)$ and fix a
Euclidean norm on the ambient space. Since $K$ is closed in $Z$ and
disjoint from $Z\cap x^\perp$, \cite[Theorem~3.11(a)]{DHR} gives,
for each $m\ge1$ and all sufficiently large $n$,
\[
s_{\alpha_n}\cdot K\subseteq U_m
   :=\{z\in E:\|z-x\|<1/m\}.
\]
Here we apply the theorem to the corresponding neighborhood in $Z$
and use the invariance of $E$. Thus $\mu(U_m)\ge\mu(K)$ for every
$m$. As $U_m\downarrow\{x\}$, continuity from above yields
$\mu(\{x\})\ge\mu(K)>0$, a contradiction.
\end{proof}

For the remaining irreducible types, a finite root system has empty
limit set, an affine root system has a single limit root, and a
rank-two indefinite root system has two limit roots
\cite[Example~2.1 and Corollary~2.15]{HLR}. In the last case the two
points form one orbit by minimality.

For a general $\Phi$, let $\Phi_i$ be its irreducible components.
By \cite[Remark~3.3]{DHR}, $E$ is the finite union of the closed
invariant sets $E(\Phi_i)$, and $W$ acts on each through $W_i$.
If an invariant probability assigned positive mass to the limit set
of an indefinite component of rank at least three, its normalized
restriction would contradict
\cref{prop:limit-root-invariant-probability}. Thus all mass lies in
the union of the singleton affine orbits and the two-point indefinite
dihedral orbits. Conversely, the uniform probability on any finite
orbit is invariant. This proves the existence criterion.

\subsection{Uniqueness of invariant measures}
\noindent\textbf{Status.} $\checkmark$ Solved. There is a unique
$W$-invariant Borel probability measure on $E$ precisely when $E$
contains exactly one finite $W$-orbit.

Indeed, all mass is carried by the finite orbits, and invariance
forces equal masses at points in the same orbit. If there is just one
such orbit, this determines the measure. Two distinct finite orbits
give distinct invariant probabilities by taking the uniform measure
on each. We count distinct orbits since component limit sets can
coincide when the simple roots are linearly dependent
\cite[Remark~3.3]{DHR}.

\subsection{Construction of invariant measures}
\noindent\textbf{Status.} $\checkmark$ Solved. Let $O_1,\dots,O_m$
be the distinct finite $W$-orbits in $E$. For $m\ge1$, all invariant
Borel probabilities are given by
\[
\boxed{\displaystyle
\mu=\sum_{j=1}^{m}\frac{t_j}{|O_j|}
                 \sum_{\xi\in O_j}\delta_\xi,\qquad
t_j\ge0,\quad \sum_{j=1}^{m}t_j=1.}
\]
Here $\delta_\xi$ is the unit point mass at $\xi$. Each group element
permutes each $O_j$, so every measure in the display is invariant.
Conversely, the weights are recovered uniquely as $t_j=\mu(O_j)$.
In particular, a singleton contributes its point mass and a two-point
orbit contributes the average of its two point masses.

\subsection{Hausdorff dimension}
Furthermore, when the root system is of indefinite
type but not hyperbolic, the set $E(\Phi)$ may exhibit
fractal properties.

This leads to another natural question:
\[
\boxed{
\text{What is the Hausdorff dimension of }E(\Phi)?
}
\]

For example, for the standard Tits realization of
the universal Coxeter group of rank four,
$E(\Phi)$ is the Apollonian gasket inscribed in a sphere,
whose Hausdorff dimension is approximately
\[
\dim_H E(\Phi)\approx 1.3057.
\]


\par\medskip
\noindent\textbf{Status.} Partially solved. Particular examples, including the Apollonian gasket, admit rigorous dimension results. No general formula for arbitrary limit-root sets was verified.

\par\medskip
\noindent\textbf{Previous work and verified progress.}
The Apollonian special case has progressed well beyond a numerical
estimate of $1.3057$. Vytnova--Wormell prove rigorous error bounds for
128 decimal places \cite[Theorem~1.1]{VW2025}; in particular,
\[
\dim_H E=1.305686728049877\ldots
\]
for that realization. Their induced transfer-operator method handles
parabolic dynamics. This does not give a formula for arbitrary Coxeter
limit-root sets, nor show that dimension depends only on the abstract
Coxeter group.

\section{Construction of Converging Sequences,
Combinatorics and Dominance Order}

Many questions concerning the precise manner in which
normalized roots converge to $E(\Phi)$ remain open.

\subsection{Convergence rates and root depth}

What convergence-rate bounds can be given in terms of root depth,
and under which hypotheses?


\par\medskip
\noindent\textbf{Status.} Partially solved. Norm--depth estimates control quadratic defect, rather than distance to a prescribed limit root.

\par\medskip
\noindent\textbf{Previous work and verified progress.}
With a linearly independent simple-root basis,
\cite[Lemma~2.10]{HLR} gives
$\|\rho\|_2^2\geq 1+c(\depth(\rho)-1)$ for some $c>0$.
Since $h(\rho)\geq\|\rho\|_2$, one obtains
\[
B(\widehat\rho,\widehat\rho)=h(\rho)^{-2}
\leq\frac{1}{1+c(\depth(\rho)-1)}.
\]
This bounds the quadratic defect, not the distance to a prescribed
limit root. Fixed-element spectral estimates are a separate special case;
a general convergence-rate theorem for arbitrary chains was not verified.

\subsection{Convergence of increasing root sequences}

Let
\[
\rho_1\leq_{\mathrm{rt}}\rho_2\leq_{\mathrm{rt}}\cdots
\]
be a sequence of positive roots increasing in the root poset.
Under what conditions does the normalized sequence
$(\widehat{\rho}_n)_{n\geq1}$ converge?

Here $\leq_{\mathrm{rt}}$ denotes the transitive closure of the
elementary steps
\[
\rho<_{\mathrm{rt}}s_i\rho
\qquad\text{whenever}\qquad
B(\rho,\alpha_i)<0.
\]
In such a step, the coefficient of $\alpha_i$ increases strictly
and all other simple-root coefficients remain unchanged.

\par\medskip
\noindent\textbf{Status.}
Partially solved. Strict increase in the root poset does not
guarantee convergence, even in the standard Tits realization
and even when consecutive roots differ by a single simple
reflection. A general criterion characterizing the convergent
increasing sequences remains open here.

\par\medskip
\noindent\textbf{Previous work and verified progress.}
By \cite[Theorem~2.7]{HLR}, every accumulation point of an
injective sequence of normalized positive roots is isotropic.
The example below shows that a single strictly increasing chain
can have accumulation points supported on disjoint subsets of
the simple roots.

\begin{proposition}
\label{prop:nonconvergent-root-poset-chain}
For every $a\geq1$, consider the rank-six based root system with
simple roots $\alpha_1,\ldots,\alpha_6$ and Gram matrix given by
\[
B(\alpha_i,\alpha_j)=
\begin{cases}
  1,  & i=j,\\
  -a, & |i-j|=1,\\
  0,  & |i-j|>1.
\end{cases}
\]
There exists a sequence of positive roots
\[
\rho_0<_{\mathrm{rt}}\rho_1<_{\mathrm{rt}}\rho_2
<_{\mathrm{rt}}\cdots
\]
in which every step is a simple reflection and whose normalized
sequence has at least two distinct accumulation points.
In particular, this occurs for the standard Tits realization
$a=1$.
\end{proposition}

\begin{proof}
Start with the positive root
\[
\rho_0=s_4\alpha_3=\alpha_3+2a\alpha_4.
\]
We interleave the following two sequences of reflections, preserving
the indicated order within each sequence:
\[
s_2,s_1,s_2,s_1,\ldots,
\qquad
s_5,s_6,s_5,s_6,\ldots.
\]
Write the current root as $\rho=\sum_{i=1}^6 c_i\alpha_i$.
Throughout the construction,
\[
c_3=1,\qquad c_4=2a.
\]
The first sequence changes only $(c_1,c_2)$, and the second
changes only $(c_5,c_6)$. Thus either sequence can be paused
and resumed without affecting its subsequent coordinate updates.

We first verify that every reflection used is an upward step.
On either side, write $(x,y)$ for the inner and outer coordinates,
respectively, and let $b$ denote the fixed central coordinate
adjacent to $x$. Thus
\[
(x,y,b)=(c_2,c_1,1)
\quad\text{or}\quad
(x,y,b)=(c_5,c_6,2a).
\]
Starting at $(x,y)=(0,0)$, the two reflections act alternately by
\[
(x,y)\longmapsto(2a(b+y)-x,y),
\qquad
(x,y)\longmapsto(x,2ax-y).
\]
Let $d_k$ be the increment of the coordinate changed at the
$k$th reflection on this side. The displayed formulas give
\[
d_1=2ab,\qquad d_2=2a\,d_1,\qquad
d_{k+1}=2a\,d_k-d_{k-1}\quad(k\geq2).
\]
Since $a\geq1$ and $b>0$, we have $d_2>d_1>0$. Moreover,
\[
d_{k+1}-d_k
=(2a-2)d_k+(d_k-d_{k-1}),
\]
so induction gives $d_{k+1}>d_k>0$ for every $k$.
At a step using $s_i$, the increment is $-2B(\rho,\alpha_i)$.
Hence every step gives $\rho<_{\mathrm{rt}}s_i\rho$.
This argument includes $a=1$, for which $d_k=2bk$.

It remains to choose the interleaving. Set
\[
L=c_1+c_2,\qquad R=c_5+c_6,\qquad C=c_3+c_4=1+2a,
\]
so that $h(\rho)=L+C+R$.
At stage $j\geq1$, apply reflections from the first sequence
until
\[
L\geq j(R+C),
\]
and then apply reflections from the second sequence until
\[
R\geq j(L+C).
\]
Each block terminates: its active coordinate sum grows without
bound if continued, while the other sum remains fixed.
Denote the roots at the ends of the two blocks by
$\rho_j^{L}$ and $\rho_j^{R}$, respectively.

Consider the linear functional
\[
F\left(\sum_{i=1}^6 z_i\alpha_i\right)=z_1+z_2.
\]
The stopping conditions give
\[
F(\widehat{\rho_j^{L}})\geq\frac{j}{j+1},
\qquad
F(\widehat{\rho_j^{R}})\leq\frac{1}{j+1}.
\]
All normalized positive roots lie in the compact simplex
$\operatorname{conv}\{\alpha_1,\ldots,\alpha_6\}$.
Thus the left and right block endpoints have accumulation
points $x_L$ and $x_R$ satisfying
\[
F(x_L)=1,\qquad F(x_R)=0.
\]
In particular, $x_L\neq x_R$. More precisely, the fixed central
coordinates vanish after normalization as the height tends to
infinity, so
\[
x_L\in\operatorname{conv}\{\alpha_1,\alpha_2\},
\qquad
x_R\in\operatorname{conv}\{\alpha_5,\alpha_6\}.
\]
Hence the full normalized sequence does not converge.
\end{proof}

\begin{remark}
For $a=2,3$, this construction uses the same Gram matrices as
the rank-six dominance-chain examples discussed below, but the
chains themselves are different. In those examples, the odd
and even roots lie in $W\alpha_3$ and $W\alpha_4$, respectively.
Indeed, for integral $a$, each simple reflection preserves every
simple-root coordinate modulo $2$, so these two orbits are
disjoint. Since comparable roots in the root poset must lie
in the same $W$-orbit, no odd term of the original dominance
chain is comparable in the root poset with an even term.
By contrast, every root in
Proposition~\ref{prop:nonconvergent-root-poset-chain}
belongs to $W\alpha_3$ and every step is explicitly upward.
\end{remark}
\subsection{Infinite reduced words and periodic root sequences}

The convergence problem for increasing root sequences leads to
the study of the possible accumulation points of
\[
\widehat{s_k s_{k-1}\cdots s_1(\alpha_0)},
\qquad k\in\mathbb N,
\]
where $\alpha_0\in\Delta$ and
\[
\cdots s_k\cdots s_1
\]
is a left-infinite reduced word of $W$.

The periodic case is of particular interest.
If the sequence of generators has period $p$, put
\[
w=s_p\cdots s_1,
\qquad
u_j=s_j\cdots s_1
\quad(0\le j<p),
\]
with $u_0=1$.
Then
\[
s_{np+j}\cdots s_1(\alpha_0)
=
u_jw^n(\alpha_0).
\]
Thus the periodic case reduces to understanding the normalized
orbits of a fixed element, together with the finitely many
intermediate phases.
For an alternating word on two generators, the corresponding
infinite dihedral subsystem provides limit roots in $E_2(\Phi)$.
Different phases may have different limits, so convergence
along complete periods does not by itself imply convergence
of the full sequence.

\par\medskip
\noindent\textbf{Status.}
Partially solved.
Theorem~\ref{thm:normalized-root-orbits} below proves that every
infinite root orbit under a fixed element has a unique normalized
limit, without a restriction on the signature of the associated
bilinear form.
Corollary~\ref{cor:periodic-root-phases} consequently determines
the accumulation set of a periodic fixed-root sequence with
an infinite period orbit.
Finite root orbits are periodic and must be treated separately.
The arguments below do not settle the convergence problem for
arbitrary nonperiodic left-infinite reduced words.

\par\medskip
\noindent\textbf{Previous work and verified progress.}
The Lorentzian spectral analysis in
\cite[Theorems~3.8 and~3.10]{ChenLabbe2017}
describes parabolic and hyperbolic iterates and identifies
their light-like eigendirections as limit roots.
For essential elements of irreducible infinite non-affine
Coxeter groups, Krammer proves that the spectral radius is
greater than one and is a simple eigenvalue
\cite[Section~6.5]{Krammer2009}.
The strict spectral gap is stated there as
Conjecture~6.5.16; it is proved below using a trace estimate,
rather than assumed.
A separate argument shows that an infinite root orbit cannot
have zero projection onto the attracting eigenspace of an
essential non-affine parabolic component.

Related results concern the inversion roots
\[
\beta_k=s_1\cdots s_{k-1}(\alpha_{s_k}).
\]
For Lorentzian systems, Chen--Labb\'e prove uniqueness of the
accumulation point of an infinite reduced-word inversion set
\cite[Theorem~2.9]{ChenLabbe2017}.
Hohlweg--Labb\'e formulate a connected-word conjecture in general
and discuss disconnected examples with several accumulation
points \cite[Section~4.2, Conjecture~3]{HL2016}.
These statements concern inversion-root sequences and should
be distinguished from the fixed-root sequences considered here.

\par\medskip
Throughout the proof, $(W,S)$ has finite rank and
\[
V=\bigoplus_{s\in S}\mathbb R\alpha_s,
\qquad
\Delta=\{\alpha_s:s\in S\},
\qquad
P=\sum_{s\in S}\mathbb R_{\geq0}\alpha_s.
\]
The symmetric bilinear form $B$ satisfies
$B(\alpha_s,\alpha_s)=1$ and the usual Coxeter root-basis
conditions; it is allowed to be degenerate.
We write
\[
h\left(\sum_{s\in S}c_s\alpha_s\right)
=\sum_{s\in S}c_s,
\qquad
\widehat{\beta}=\frac{\beta}{h(\beta)}
\quad(\beta\in\Phi).
\]
Let $\mathcal T\subseteq V^*$ denote the dual Tits cone.
An element $u\in W$ is called \emph{essential} if its parabolic
closure satisfies $\operatorname{Pc}(u)=W$.

\begin{lemma}
\label{lem:periodic-trace-bound}
For every $g\in W$,
\[
\operatorname{tr}_V(g)\geq \dim V-2\ell(g).
\]
Consequently,
\[
\operatorname{tr}_V(g^n)
\geq \dim V-2n\ell(g)
\qquad(n\geq1).
\]
\end{lemma}

\begin{proof}
We first show that
\[
\ell(vs)=\ell(sv)=\ell(v)+1
\quad\Longrightarrow\quad
B(\alpha_s,v\alpha_s)\leq1.
\]
The length assumptions imply that both $v\alpha_s$ and
$v^{-1}\alpha_s$ are positive roots.
Suppose that $B(\alpha_s,v\alpha_s)>1$.
The dominance criterion
\cite[Proposition~1.4.7(b)]{Krammer2009}
implies that $\alpha_s$ and $v\alpha_s$ are comparable.
A simple root cannot strictly dominate another positive root:
indeed, its simple reflection makes it negative while preserving
every other positive root.
It follows that
\[
v\alpha_s\succ\alpha_s.
\]
By $W$-invariance of dominance,
\[
\alpha_s\succ v^{-1}\alpha_s,
\]
which is again impossible.
This proves the claimed inequality.

We now use strong induction on $\ell(g)$.
The assertion is immediate for $g=1$.
Choose an element $z$ of minimal length in the conjugacy class
of $g$.
If $z\neq1$, write
\[
z=vs,
\qquad
\ell(v)=\ell(z)-1.
\]
Since $sv=szs$ is conjugate to $z$, minimality gives
\[
\ell(sv)=\ell(vs)=\ell(v)+1.
\]
The reflection formula yields
\[
\operatorname{tr}_V(vs)
=
\operatorname{tr}_V(v)-2B(\alpha_s,v\alpha_s).
\]
Therefore, by the preceding inequality and the induction
hypothesis,
\[
\begin{aligned}
\operatorname{tr}_V(g)
&=\operatorname{tr}_V(z)\\
&\geq\operatorname{tr}_V(v)-2\\
&\geq\dim V-2\ell(z)\\
&\geq\dim V-2\ell(g).
\end{aligned}
\]
The assertion for powers follows from
$\ell(g^n)\leq n\ell(g)$.
\end{proof}

\begin{lemma}
\label{lem:periodic-strict-gap}
Suppose that $W$ is irreducible, infinite and non-affine,
and let $u\in W$ be essential.
Then the spectral radius $\rho$ of $u$ satisfies $\rho>1$,
is an algebraically simple eigenvalue, and
\[
|\lambda|<\rho
\]
for every other complex eigenvalue $\lambda$ of $u$.
\end{lemma}

\begin{proof}
Set
\[
R=\operatorname{rad}B,
\qquad
\overline V=V/R,
\]
and let $A$ denote the action induced by $u$ on $\overline V$.
The quotient is again a root-basis realization of $W$
by \cite[Proposition~6.1.3]{Krammer2009}.

Krammer's results
\cite[Lemma~6.5.4 and Theorem~6.5.14]{Krammer2009}
show that the spectral radius $\rho>1$ of $A$ is an
algebraically simple eigenvalue and that all eigenvalues
of modulus $\rho$ are semisimple.
Moreover, every positive power of $u$ is essential.
Indeed, the set of odd reflections is unchanged on passing
to a positive power, and these reflections generate the
infinite components of the parabolic closure
\cite[Lemma~5.2.7 and Corollary~5.8.7]{Krammer2009}.
Hence $\rho^m$ is an algebraically simple eigenvalue of $A^m$
for every $m\geq1$.

Every element of $W$ fixes $R$ pointwise.
Lemma~\ref{lem:periodic-trace-bound} therefore gives
\[
\operatorname{tr}_{\overline V}(A^n)
\geq\dim\overline V-2n\ell(u).
\]
Let $F\subseteq\overline V$ be the real spectral subspace
corresponding to all eigenvalues of modulus $\rho$, and put
\[
T=\rho^{-1}A|_F.
\]
Since $T$ is semisimple and all its eigenvalues have modulus
one, the closure
\[
K=\overline{\{T^n:n\in\mathbb Z\}}
\]
is a compact abelian group.
Define
\[
f(k)=\operatorname{tr}_F(k)
\qquad(k\in K).
\]
This function is real-valued, and
\[
\rho^{-n}\operatorname{tr}_{\overline V}(A^n)
=
f(T^n)+o(1).
\]
Every element of $K$ is a limit of a sequence $T^{n_j}$ with
$n_j\to\infty$.
The trace lower bound consequently implies
\[
f(k)\geq0
\qquad(k\in K).
\]

Let $\mu$ be normalized Haar measure on $K$.
Decompose the complex representation $F_{\mathbb C}$ into
characters of $K$:
\[
f=\sum_{\chi}m_\chi\chi.
\]
The multiplicity of the trivial character is one, since
the eigenvalue $\rho$ of $A$ is simple.
Thus
\[
\int_K f\,d\mu=1.
\]

Suppose that a nontrivial character $\chi$ occurs.
By character orthogonality,
\[
m_\chi=\int_K f\,\overline{\chi}\,d\mu.
\]
Since $f\geq0$ and $|\chi|=1$,
\[
1\leq m_\chi
\leq\int_K f\,d\mu=1.
\]
It follows that $m_\chi=1$ and
\[
\int_K f(k)\bigl(1-\operatorname{Re}\chi(k)\bigr)\,d\mu(k)=0.
\]
The integrand is continuous and nonnegative.
Hence $\chi(k)=1$ whenever $f(k)>0$.
Since
\[
f(1)=\dim F>0,
\]
the kernel of $\chi$ contains a neighborhood of the identity.
It is therefore an open subgroup of the compact group $K$,
so $\chi(K)$ is finite.

Choose $m\geq1$ such that $\chi(T)^m=1$.
The trivial-character space and the $\chi$-space then both
contribute to the eigenvalue $\rho^m$ of $A^m$.
This contradicts the algebraic simplicity of that eigenvalue.
Thus no nontrivial character occurs, proving that $\rho$ is
the unique eigenvalue of modulus $\rho$.

Finally, the action on $R$ is the identity.
Passing from $\overline V$ back to $V$ only adds eigenvalues
equal to $1$, so the conclusion also holds on $V$.
\end{proof}

\begin{lemma}
\label{lem:periodic-ambient-root}
Let $J\subseteq S$ be such that $W_J$ is irreducible,
infinite and non-affine, and let $u\in W_J$ satisfy
\[
\operatorname{Pc}_{W_J}(u)=W_J.
\]
Let $\rho>1$ be its spectral radius on $V_J$.
For every root $\beta\in\Phi$ with infinite $u$-orbit,
there exists $v\in V_J\setminus\{0\}$ such that
\[
\rho^{-n}u^n\beta\longrightarrow v.
\]
\end{lemma}

\begin{proof}
Let $\mathcal T_J\subseteq V_J^*$ denote the dual Tits cone
of the subsystem.
The axis results
\cite[Theorems~6.5.10 and~6.5.14]{Krammer2009}
provide nonzero linear functionals $f_+,f_-\in V_J^*$ such that,
after rescaling,
\[
f_\pm\circ u=\rho^{\pm1}f_\pm,
\qquad
x_J:=f_++f_-\in\operatorname{int}\mathcal T_J.
\]
If $B|_{V_J}$ is degenerate, apply these results first to
$V_J/\operatorname{rad}(B|_{V_J})$ and pull the functionals
back to $V_J$.
The pulled-back axis point is still an interior point:
after moving it into the fundamental closed chamber, its
zero set of simple roots is spherical, exactly as in the
quotient realization.
We use here the characterization of the interior of the
Tits cone by facets with finite stabilizer
\cite[Corollary~2.2.5]{Krammer2009}.

Since
\[
(u-I)V\subseteq V_J,
\]
define extensions to $V$ by
\[
F_\pm(z)
=
\frac{f_\pm((u-I)z)}{\rho^{\pm1}-1}.
\]
They satisfy
\[
F_\pm|_{V_J}=f_\pm,
\qquad
F_\pm\circ u=\rho^{\pm1}F_\pm.
\]

There exists $X\in\operatorname{int}\mathcal T$ whose
restriction to $V_J$ is $x_J$.
To see this, move $x_J$ into a fundamental face for $W_J$.
Its zero set is spherical.
Extend it by assigning strictly positive values to all
simple roots outside $J$, and then move the extension back.
The resulting functional lies in the interior of the
ambient Tits cone by the same finite-stabilizer criterion.

Consequently,
\[
X=F_++F_-+\eta,
\qquad
\eta|_{V_J}=0.
\]
Since $(u-I)V\subseteq V_J$, we have $\eta\circ u=\eta$, and
therefore
\[
X(u^n\beta)
=
\rho^nF_+(\beta)
+\rho^{-n}F_-(\beta)
+\eta(\beta).
\]

For every $X\in\operatorname{int}\mathcal T$ and every
$M>0$, only finitely many roots $\gamma$ satisfy
$|X(\gamma)|\leq M$.
This follows by applying
\cite[Lemma~5.7.1]{Krammer2009}
to $\gamma$ and $-\gamma$.
Since an infinite cyclic root orbit has no repetitions,
\[
|X(u^n\beta)|\longrightarrow\infty.
\]
The preceding formula then forces
\[
F_+(\beta)\neq0.
\]

By Lemma~\ref{lem:periodic-strict-gap}, $\rho$ is simple
and strictly dominant on $V_J$.
The action induced by $u$ on $V/V_J$ is the identity,
so the same assertion holds on $V$.
Let $P_\rho$ be the corresponding spectral projection.
Because $F_+\circ u=\rho F_+$,
\[
F_+=F_+\circ P_\rho.
\]
Thus $P_\rho\beta\neq0$, and the strict spectral gap gives
\[
\rho^{-n}u^n\beta\longrightarrow P_\rho\beta.
\]
Finally, a $\rho$-eigenvector belongs to $V_J$, since $u$
acts trivially on $V/V_J$ and $\rho\neq1$.
\end{proof}

\begin{theorem}[Convergence of normalized root orbits]
\label{thm:normalized-root-orbits}
Let $w\in W$ and $\alpha\in\Phi$.
If
\[
\mathcal O_w(\alpha)
:=
\{w^n\alpha:n\geq0\}
\]
is infinite, then there exists a unique point
$x_{w,\alpha}\in E(\Phi)$ such that
\[
\widehat{w^n\alpha}\longrightarrow x_{w,\alpha}.
\]

If $\mathcal O_w(\alpha)$ is finite, the normalized sequence
is periodic and converges if and only if
$w\alpha=\pm\alpha$, in which case it is constant with
value $\widehat{\alpha}$.
\end{theorem}

\begin{proof}
We first establish the required asymptotics for affine
parabolic components.

Suppose that $W_J$ is irreducible affine and that
$B|_{V_J}$ is positive semidefinite, with
\[
\operatorname{rad}(B|_{V_J})=\mathbb R\delta,
\qquad
\delta\in P\setminus\{0\}.
\]
Every element of $W_J$ fixes $\delta$.
For $u\in W_J$, choose $q\geq1$ such that $u^q$ acts
trivially on $V_J/\mathbb R\delta$.
Such a $q$ exists because the action on this quotient has
finite image.
Set
\[
N=u^q-I.
\]
On the ambient space $V$,
\[
NV\subseteq V_J,
\qquad
NV_J\subseteq\mathbb R\delta,
\qquad
N\delta=0.
\]
Hence $N^3=0$, and for every $\beta\in\Phi$,
\[
u^{qn}\beta
=
\beta+nN\beta+\binom n2N^2\beta.
\]

If $N^2\beta\neq0$, then $N^2\beta\in\mathbb R\delta$ is
fixed by $u$.
For $0\leq r<q$,
\[
u^{qn+r}\beta
=
u^r\beta+n\,u^rN\beta+\binom n2N^2\beta.
\]
It follows, uniformly over the finitely many residues $r$,
that
\[
u^m\beta
=
m^2\frac{N^2\beta}{2q^2}+O(m).
\]

If $N^2\beta=0$ but $N\beta\neq0$, invariance of $B$ gives
\[
B(\beta+nN\beta,\beta+nN\beta)=B(\beta,\beta)=1.
\]
Comparing the quadratic coefficients yields
\[
B(N\beta,N\beta)=0.
\]
Since $N\beta\in V_J$ and $B|_{V_J}$ is positive
semidefinite,
\[
N\beta\in\mathbb R\delta.
\]
Thus $uN\beta=N\beta$, and
\[
u^m\beta=m\frac{N\beta}{q}+O(1).
\]

If $N\beta=0$, then $u^q\beta=\beta$, so the orbit is finite.
We have therefore proved that every infinite orbit of this
affine component satisfies
\[
u^n\beta=n^d v+O(n^{d-1}),
\qquad
d\in\{1,2\},
\qquad
0\neq v\in\mathbb R\delta.
\]

There is one additional realization of an abstract affine
component to consider.
If $W_J=\langle s,t\rangle$ is infinite dihedral and
\[
B(\alpha_s,\alpha_t)=-c,
\qquad c>1,
\]
then $B|_{V_J}$ is nondegenerate.
An essential element is $(st)^k$ for some $k\neq0$, and
its eigenvalues on $V_J$ are
\[
\rho=(c+\sqrt{c^2-1})^{2|k|}>1,
\qquad
\rho^{-1}.
\]
Moreover,
\[
V=V_J\oplus V_J^\perp,
\]
and $W_J$ fixes $V_J^\perp$ pointwise.
If the $\rho$-component of $\beta$ were zero, its forward
orbit would be bounded and hence finite, since the root
system is locally finite.
Thus every infinite orbit in this case satisfies
\[
\rho^{-n}u^n\beta\longrightarrow v
\qquad
\text{for some }v\in V_J\setminus\{0\}.
\]
All irreducible affine components of rank at least three
have the usual positive semidefinite Gram form, so the
preceding cases exhaust the affine possibilities.

We now treat arbitrary $w$.
First suppose that its parabolic closure is standard:
\[
\operatorname{Pc}(w)=W_I.
\]
Write
\[
I=I_1\sqcup\cdots\sqcup I_k
\]
for the decomposition into irreducible components.
Then
\[
w=u_1\cdots u_k,
\qquad
u_i\in W_{I_i},
\qquad
\operatorname{Pc}_{W_{I_i}}(u_i)=W_{I_i}.
\]
The factors commute, and each factor fixes the root spaces
of the other components pointwise.
In particular,
\[
(u_i^n-I)V\subseteq V_{I_i},
\qquad
(u_j^n-I)V_{I_i}=0
\quad(i\neq j).
\]
Expanding the product gives the exact identity
\[
w^n\alpha
=
\alpha+\sum_{i=1}^k(u_i^n\alpha-\alpha).
\]

If every $u_i$-orbit of $\alpha$ were finite, a common
multiple of their periods would also be a period for the
$w$-orbit.
Hence at least one component orbit is infinite.

For a non-affine infinite component,
Lemma~\ref{lem:periodic-ambient-root} shows that an infinite
component orbit satisfies
\[
u_i^n\alpha=\rho_i^n v_i+o(\rho_i^n),
\qquad
\rho_i>1,
\qquad
0\neq v_i\in V_{I_i}.
\]
The same form of asymptotics holds for an infinite
dihedral component with parameter $c>1$.
An infinite orbit in a positive semidefinite affine
component has the polynomial asymptotics established above.
All remaining component orbits are finite and therefore
bounded.

Suppose first that some component orbit has exponential
growth, and let
\[
R=\max\{\rho_i:
        \text{the $u_i$-orbit of $\alpha$ is infinite
        and has exponential growth}\}.
\]
Then
\[
R^{-n}w^n\alpha
\longrightarrow
v:=\sum_{\rho_i=R}v_i.
\]
The nonzero summands belong to distinct coordinate subspaces
$V_{I_i}$, so $v\neq0$.

If no component orbit has exponential growth, let
$d\in\{1,2\}$ be the largest degree of an infinite affine
component orbit.
Then
\[
n^{-d}w^n\alpha
\longrightarrow
v:=\sum_{d_i=d}v_i\neq0.
\]
Thus in either case there is a sequence of positive numbers
$a_n\to\infty$ such that
\[
a_n^{-1}w^n\alpha\longrightarrow v\neq0.
\]

For general $w$, choose $g\in W$ such that
\[
p:=gwg^{-1}
\]
has standard parabolic closure, and set $\beta=g\alpha$.
The orbit of $\beta$ under $p$ is infinite.
Applying the preceding argument gives
\[
a_n^{-1}p^n\beta\longrightarrow v\neq0.
\]
Consequently,
\[
a_n^{-1}w^n\alpha
=
g^{-1}\bigl(a_n^{-1}p^n\beta\bigr)
\longrightarrow g^{-1}v\neq0.
\]
We may therefore assume, without restriction, that
\[
a_n^{-1}w^n\alpha\longrightarrow v\neq0
\]
for the original $w$ and $\alpha$.

Every root belongs to $P\cup(-P)$.
Since $a_n>0$ and this union is closed,
\[
v\in(P\cup(-P))\setminus\{0\}.
\]
Hence $h(v)\neq0$, and
\[
\widehat{w^n\alpha}
=
\frac{a_n^{-1}w^n\alpha}
     {h(a_n^{-1}w^n\alpha)}
\longrightarrow
\frac{v}{h(v)}.
\]

The roots in an infinite cyclic orbit are pairwise distinct.
No two can be negatives of one another, since
$w^m\alpha=-\alpha$ would imply $w^{2m}\alpha=\alpha$.
Because all roots have squared length one, equal normalized
roots differ only by a sign.
Thus the normalized roots in the orbit are also pairwise
distinct.
Their limit therefore belongs to $E(\Phi)$.

Finally, if the root orbit is finite, it is periodic.
A periodic sequence converges precisely when it is constant.
Here constancy is equivalent to
$\widehat{w\alpha}=\widehat{\alpha}$, which, by preservation
of squared length, is equivalent to $w\alpha=\pm\alpha$.
\end{proof}

\begin{remark}
The assumption in
Theorem~\ref{thm:normalized-root-orbits}
concerns the infinitude of the root orbit, not merely
the order of the group element.
An infinite-order element may have finite root orbits.

The conclusion also extends to a finite based root system
whose simple roots are only positively independent.
Choose a linear functional $h$ strictly positive on the
simple roots, and lift the realization to the vector space
having those roots as a formal basis.
The natural map to the original realization is equivariant
and cannot annihilate a nonzero vector in the positive cone.
Applying the preceding proof to the lifted root orbit,
with the pulled-back normalizing functional, and then
projecting gives the same convergence conclusion.
\end{remark}

\begin{corollary}[Intermediate phases of a periodic word]
\label{cor:periodic-root-phases}
Suppose that the generator sequence has period $p$, and put
\[
w=s_p\cdots s_1,
\qquad
u_j=s_j\cdots s_1
\quad(0\leq j<p),
\qquad
u_0=1.
\]
Assume that the orbit $\{w^n\alpha_0:n\geq0\}$ is infinite,
and let
\[
x=\lim_{n\to\infty}\widehat{w^n\alpha_0}.
\]
Then every phase converges:
\[
\lim_{n\to\infty}
\widehat{s_{np+j}\cdots s_1(\alpha_0)}
=
u_j\cdot x
=
\widehat{u_jx}.
\]
The accumulation set of the full normalized sequence is
exactly
\[
\{u_j\cdot x:0\leq j<p\}.
\]
In particular, the full sequence converges if and only if
\[
u_j\cdot x=x
\qquad(0\leq j<p).
\]
\end{corollary}

\begin{proof}
By the proof of
Theorem~\ref{thm:normalized-root-orbits},
there exist positive numbers $a_n\to\infty$ and a nonzero
vector $v$ such that
\[
a_n^{-1}w^n\alpha_0\longrightarrow v,
\qquad
x=\widehat v.
\]
For each $j$,
\[
a_n^{-1}u_jw^n\alpha_0\longrightarrow u_jv\neq0.
\]
The vectors $u_jw^n\alpha_0$ are roots.
Thus
\[
u_jv\in(P\cup(-P))\setminus\{0\},
\]
so $h(u_jv)\neq0$, and normalization gives
\[
\widehat{u_jw^n\alpha_0}
\longrightarrow
\widehat{u_jv}
=
u_j\cdot x.
\]
Each phase has infinitely many terms, so each displayed
limit is an accumulation point.
Conversely, every convergent subsequence of the full
sequence has a further subsequence with a fixed residue
class modulo $p$, and its limit must be the corresponding
phase limit.
The assertion about convergence follows.
\end{proof}

\begin{remark}
If the period orbit of $\alpha_0$ is finite, choose $q\geq1$
with $w^q\alpha_0=\alpha_0$.
Then the full root sequence is periodic with period dividing
$pq$.
Its normalized sequence converges exactly when all terms
in one such period have the same normalization.
Thus finite period orbits are also completely determined,
but their accumulation points need not be limit roots.
\end{remark}
\subsection{Uniqueness of limits of dominance chains}

Let $(\rho_n)_{n\ge1}$ be an infinite sequence of positive roots
which is strictly increasing in dominance order:
\[
\rho_1\prec\rho_2\prec\cdots.
\]
Must the normalized sequence $(\widehat{\rho}_n)_{n\ge1}$
converge to a unique point of $E(\Phi)$?

\par\medskip
\noindent\textbf{Status.}
\(\checkmark\) \textit{Answered negatively in general in
\cite{WangUniqueness}.}
That manuscript proves uniqueness in rank at most five and
for nondegenerate systems of signature $(r-1,1)$ or $(2,r-2)$.
It also constructs rank-six strict dominance chains with exactly
two distinct normalized accumulation points, in signatures
$(3,3)$ and $(4,2)$.
These counterexamples use Gram parameters strictly below $-1$
on infinite edges; they do not settle the general question
restricted to the standard Tits realization, where those
parameters are fixed at $-1$.

\par\medskip
\noindent\textbf{Previous work and subsequent progress.}
Any two accumulation points of a normalized strict dominance
chain are isotropic and mutually orthogonal
\cite[Proposition~6.8(b)]{DHR}.
This gives uniqueness in the Lorentzian case, but orthogonality
alone does not imply proportionality when the Witt index
is at least two.

Fu--Reeves--Xu prove uniqueness when a chain has an accumulation
point arising from an affine reflection subgroup
\cite[Proposition~6.5]{FRX2026}.
The preprint \cite{FuEtAl2026} establishes rank-three uniqueness
(Proposition~4.5), identifies the accumulation supports with
the eventual chain support (Theorem~4.8), and proves uniqueness
when an accumulation point is Pos-finite (Proposition~4.14).
Its announcement of a forthcoming general treatment should
be distinguished from the results proved there.

The supplied manuscript \cite{WangUniqueness} provides proofs
of further positive cases and explicit higher-rank
counterexamples, as follows.

\begin{theorem}\label{thm:main}
For the geometric realizations considered here, allowing arbitrary
Gram parameters at most $-1$ on infinite edges, the following
statements hold.
\begin{enumerate}
    \item If $|S|\le5$, every infinite strict dominance chain
    has a normalized sequence converging to a unique point
    of $E(\Phi)$.

    \item In arbitrary rank $r$, the same conclusion holds
    for every nondegenerate system of signature
    \[
    (r-1,1)
    \qquad\text{or}\qquad
    (2,r-2).
    \]

    \item More generally, if the system is irreducible and
    its inertia is $(p,1,k)$, with $p,k\ge1$, every infinite
    strict dominance chain has a unique normalized limit.

    \item There exist irreducible rank-six systems of
    signatures $(3,3)$ and $(4,2)$ admitting infinite strict
    dominance chains whose normalized sequences have exactly
    two distinct accumulation points.
\end{enumerate}
\end{theorem}

\begin{proof}[Sources and proof structure]
Parts~(1) and~(2) are
\cite[Theorem~2.1]{WangUniqueness}.
The proof combines the small-rank inertia analysis, quotient
lifting, the relative-boundary argument in signature $(3,2)$,
and its extension to positive index two.
Part~(3) is
\cite[Corollary~2.18]{WangUniqueness}.
Part~(4) is
\cite[Proposition~4.10 and Remark~4.11]{WangUniqueness}.
\end{proof}

\par\medskip
\noindent\textbf{The rank-six counterexamples.}
For $a>1$, consider the path Gram matrix
\[
G_a=
\begin{pmatrix}
1&-a&0&0&0&0\\
-a&1&-a&0&0&0\\
0&-a&1&-a&0&0\\
0&0&-a&1&-a&0\\
0&0&0&-a&1&-a\\
0&0&0&0&-a&1
\end{pmatrix}.
\]
Its eigenvalues are
\[
1-2a\cos\frac{k\pi}{7},
\qquad 1\le k\le6.
\]
Thus $G_3$ has signature $(3,3)$ and $G_2$ has signature $(4,2)$.

Here is the structure of the construction in
\cite[Section~4.3]{WangUniqueness}.
Put
\[
g=s_1s_2,
\qquad
t=s_5s_6,
\qquad
r=a+\sqrt{a^2-1},
\]
and define
\[
x=\frac{r\alpha_1+\alpha_2}{r+1},
\qquad
y=\frac{r\alpha_5+\alpha_6}{r+1}.
\]
These are distinct orthogonal limit roots.
The positive roots
\[
u_m=g^m\alpha_3,
\qquad
v_m=t^m\alpha_4
\]
satisfy
\[
\widehat{u_m}\longrightarrow x,
\qquad
\widehat{v_m}\longrightarrow y,
\qquad
B(u_m,v_\ell)=-a
\quad(m,\ell\ge0).
\]

Choose rapidly increasing integers $M_n$, and put
\[
\gamma_n=
\begin{cases}
u_{M_n},&n\text{ odd},\\
v_{M_n},&n\text{ even},
\end{cases}
\qquad
w_0=1,
\qquad
w_n=s_{\gamma_1}\cdots s_{\gamma_n},
\qquad
\beta_n=w_{n-1}\gamma_n.
\]
The constant cross-pairing gives
\[
s_{\gamma_n}\gamma_{n+1}
=\gamma_{n+1}+2a\gamma_n
\succ\gamma_n.
\]
By $W$-invariance of dominance,
\[
\beta_{n+1}\succ\beta_n.
\]
The roots $\beta_n$ are positive, so this is an infinite strict
dominance chain.

The integers $M_n$ are chosen so that
$\widehat{\beta_n}$ alternately approximates the moving endpoints
\[
\widehat{w_{n-1}x},
\qquad
\widehat{w_{n-1}y},
\]
with summable errors.
On each odd step, the second endpoint is fixed; on each even
step, the first endpoint is fixed.
The resulting endpoint sequences converge to distinct points
$X,Y\in E(\Phi)$, and
\[
\widehat{\beta}_{2n-1}\longrightarrow X,
\qquad
\widehat{\beta}_{2n}\longrightarrow Y,
\qquad
X\ne Y,
\qquad
B(X,Y)=0.
\]
For $a=3$, the manuscript gives the explicit choice $M_n=10^n$.
For $a=2$, it chooses the exponents recursively.
See \cite[Proposition~4.10 and Remark~4.11]{WangUniqueness}
for the estimates and the complete proof.

\begin{remark}
Adjoining $r-6$ orthogonal finite rank-one factors preserves
these chains and yields counterexamples of signatures
\[
(r-3,3)
\qquad\text{and}\qquad
(r-2,2)
\qquad(r\ge6).
\]
These higher-rank extensions are reducible; the rank-six
examples themselves are irreducible.

The constructions are existence results.
They do not assert failure for every system outside the
positive families above, nor do they provide counterexamples
in every remaining signature.
A complete classification by signature is not established.
\end{remark}

\begin{remark}
The counterexamples here concern actual strict dominance chains.
They therefore give a negative answer to general uniqueness
independently of the counterexamples concerning convex hulls,
facial restriction, or exposed points.

However, the path constructions use $a=2$ or $a=3$.
Setting $a=1$ makes the two outer dihedral blocks affine and
invalidates the exponential-attraction argument used above.
The general uniqueness question for standard Tits realizations
is consequently not resolved by these examples.
\end{remark}
\subsection{Ends of the dominance order}

One should distinguish limits of dominance-increasing chains from
limits of auxiliary sequences associated with such chains; see
\cite[Proposition~6.8(a) and Section~7.7]{DHR}. The latter statement
alone does not settle which points are actual chain limits.

These observations suggest associating subsets of
$E(\Phi)$ with the ends of the dominance order.


\par\medskip
\noindent\textbf{Status.} No general solution verified. Existing chain constraints do not construct a complete boundary of the dominance poset.

\par\medskip
\noindent\textbf{Previous work and verified progress.}
Fu's geometric interpretation of dominance via the imaginary cone
\cite{Fu2013} is a relevant starting point. The support and sign constraints
of \cite[Theorems~4.8 and~4.18]{FuEtAl2026} further restrict the accumulation
sets of chains. They do not, as stated, construct a complete boundary of
the dominance poset. In particular, minimality of the action cannot
replace an argument that each point is an actual chain limit.

\subsection{Ends of the weak order and isotropic faces}

Similarly, the results of Dyer suggest a possible
approach to associating isotropic faces of the imaginary
cone with the ends of the weak order.

Although these ideas have been studied most extensively
for generic universal root systems, several basic
questions remain open even in that setting.

Clarifying these constructions and their relationships
would improve our understanding of the connection
between the combinatorics of root systems and the
distribution of normalized roots.


\par\medskip
\noindent\textbf{Status.} No general solution verified. Results on inversion sets do not establish the full proposed correspondence with isotropic faces.

\par\medskip
\noindent\textbf{Previous work and verified progress.}
Hohlweg--Labb\'e develop the link between infinite reduced words,
biclosed sets, and paths in the imaginary cone
\cite[Sections~4.2--4.3]{HL2016}. Their connected-word and path convergence
statements are formulated as conjectures in the general setting.
The Lorentzian inversion-set theorem cited above resolves an important
special case, but does not itself produce the proposed correspondence
with all isotropic faces in arbitrary signature.

\subsection{Combinatorial determination of the closed imaginary cone}

Finally, the Coxeter system $(W,S)$, together with
its Coxeter graph and the collection of facial subsets
of $\Delta$, determines the face lattice of the
imaginary cone as a lattice equipped with a $W$-action.

It is natural to ask whether adding the information
of the affine edges
\[
\left\{
\{\alpha,\beta\}\subseteq\Delta
\mid B(\alpha,\beta)=-1
\right\}
\]
is sufficient to determine the face lattice of the
closed imaginary cone combinatorially.


\par\medskip
\noindent\textbf{Status.} Partially solved. Face-lattice results for the imaginary cone do not automatically describe all faces of its topological closure.

\par\medskip
\noindent\textbf{Previous work and verified progress.}
Dyer describes the face lattice of the imaginary cone in terms of
facial reflection subgroups \cite{Dyer2019}. This concerns the cone
before taking its topological closure. Additional faces of the closure
are the issue in the present question; a classification for the unclosed
cone cannot be cited as an answer for the closed cone. No general
resolution from precisely the listed combinatorial data was verified.

\subsection{Combinatorial determination of the limit-root set}

It is also conceivable that the limit-root set admits
a combinatorial description in terms of this data,
determining it either as a set or at least up to
homeomorphism.


\par\medskip
\noindent\textbf{Status.} No general solution verified. Reconstruction of an embedded set and reconstruction up to homeomorphism must be distinguished.

\par\medskip
\noindent\textbf{Previous work and verified progress.}
In favorable representations, boundary models give combinatorial
access to the topology; see the conditional boundary discussion above.
The deformation theory in \cite[Section~8]{DGKLM2025} also shows why
representation parameters must be tracked. Neither it nor the face-lattice
result in \cite{Dyer2019} asserts that the stated data determine every
embedded limit-root set. A set-theoretic reconstruction and a classification
up to equivariant homeomorphism should be posed separately.

\section{Generalization to Other Frameworks}

\subsection{Extensions to other root-system frameworks}

The concept of root systems appears in several
different mathematical frameworks, including:
\begin{enumerate}
    \item Coxeter groups;
    \item semisimple Lie algebras;
    \item Kac--Moody Lie algebras;
    \item extended affine Lie algebras;
    \item reductive algebraic groups.
\end{enumerate}

In most of these settings, limit roots can still be
defined, and in some cases an isotropic cone can
also be defined.

It is expected that some of the results established
in the present work and in \cite{HLR} can be generalized
to these other settings.


\par\medskip
\noindent\textbf{Status.} Research direction with established partial results. Kac--Moody and representation-theoretic connections do not constitute a single extension theorem for all listed frameworks.

\par\medskip
\noindent\textbf{Previous work and verified progress.}
Dyer's imaginary-cone construction connects general Coxeter root
systems to the Kac--Moody setting \cite{Dyer2019}. Paquette provides a
concrete representation-theoretic development for quivers
\cite{Paquette2016}. These are substantial bridges, not a single extension
theorem covering all five frameworks. In finite root systems there is
no infinite set of root directions to accumulate, so a nontrivial limit-root
question requires an infinite setting or an additional asymptotic construction.

\subsection{Limit roots in quiver representation theory}

For example, certain classes of based root systems
arise naturally in the theory of quiver representations,
where positive roots can be interpreted as dimension
vectors of indecomposable representations.

This leads to the following question:

\[
\boxed{
\begin{gathered}
\text{How do limit roots arise in}\\
\text{quiver representation theory?}
\end{gathered}
}
\]


\par\medskip
\noindent\textbf{Status.} Partially solved. Paquette describes rational accumulation points of real Schur roots and irrational accumulation points in the weakly hyperbolic case. This does not classify all real roots of arbitrary quivers.

\par\medskip
\noindent\textbf{Previous work and verified progress.}
Paquette studies projective accumulation of real Schur roots,
which encode exceptional representations of an acyclic quiver
\cite{Paquette2016}. Canonical decomposition describes the rational
accumulation points, and the weakly hyperbolic case admits a complete
description of the irrational ones. This directly develops the direction
suggested in the question. The distinction between real Schur roots and
all real roots is essential; the results should not be stated as a
classification for every Kac--Moody or Coxeter root system.

\bibliographystyle{unsrt}
\bibliography{reference}
\end{document}